% !TeX program = lualatex
% ================================================================
% Punjabi Mathematical Dictionary
%
% Generated by the server using the following settings:
%
%   language            Punjabi (pan)
%   text direction      left to right
%   font                Noto Sans Gurmukhi
%   fallback font       Noto Serif
%   built from          latex/dictionary/mathematicalDictionary.tex
%   tier 3 vocabulary   green   RGB 0 128 0
%   English text        black   RGB 0 0 0
%   Punjabi text        purple  RGB 112 48 160
%   layout macros       version 3
%
% Please don't edit any information in this comment block - the server
% compares it against the current settings and rebuilds this file when
% they differ, so changes here would be overwritten. However, feel free
% to make updates to the LaTeX below if you want to edit it.
% ================================================================
%
% ---- What is safe to edit in this file ----------------------------------
%
% Safe:     the Punjabi word and the Punjabi definition - the third
%           and fourth arguments of each entry. That is what this file is for:
%           if a translation reads wrongly to somebody who speaks the language,
%           correct it here and every sheet that uses the word follows.
%
% Not safe: the English word and the English definition - the first and second
%           arguments. The first is how an entry is matched to the shared
%           dictionary; the second is the wording the translation was made
%           from, and is how the server knows the translation is still current.
%           Changing an English definition here does not change the shared
%           dictionary - it only makes this entry look up to date when it is
%           not. Edit the English in the shared dictionary instead, and this
%           file will be brought back into step on the next run.
%
% Entries are added and removed by the server, following the shared dictionary.
% An entry the server cannot read is reported rather than guessed at, so if a
% run complains about this file, look for a missing brace.
% -------------------------------------------------------------------------

\documentclass[11pt]{article}
\usepackage[a4paper,margin=18mm]{geometry}
% Characters this language's font has not got are borrowed from another,
% rather than being silently dropped - see docs/translated-sheet-latex.md
\directlua{%
  if luaotfload and luaotfload.add_fallback then
    luaotfload.add_fallback("eallatinfallback",
      { "Noto Serif:mode=harf;" })
  else
    tex.error("this luaotfload is too old to register a fallback font")
  end
}
\usepackage[english]{babel}
\babelprovide[import]{punjabi}
\babelfont[punjabi]{rm}[Renderer=Harfbuzz, BoldFont={Noto Sans Gurmukhi}, ItalicFont={Noto Sans Gurmukhi}, BoldItalicFont={Noto Sans Gurmukhi}, RawFeature={fallback=eallatinfallback}]{Noto Sans Gurmukhi}
\babelfont[punjabi]{sf}[Renderer=Harfbuzz, BoldFont={Noto Sans Gurmukhi}, ItalicFont={Noto Sans Gurmukhi}, BoldItalicFont={Noto Sans Gurmukhi}, RawFeature={fallback=eallatinfallback}]{Noto Sans Gurmukhi}
\newcommand{\ealtext}[1]{\foreignlanguage{punjabi}{#1}}
\newcommand{\ealtextblock}[1]{{\raggedright\foreignlanguage{punjabi}{#1}\par}}
% ================================================================
% EAL layout helpers
% ================================================================
\usepackage{xcolor}

\definecolor{ealkeycolour}{RGB}{0,128,0}
\definecolor{ealenglish}{RGB}{0,0,0}
\definecolor{ealtwocolour}{RGB}{112,48,160}

\newsavebox{\ealboxen}
\newsavebox{\ealboxtr}
\newlength{\ealglosswd}
\newlength{\ealglossraise}
\setlength{\ealglossraise}{1.15em}

% a tier 3 word, in English and in translation alike
\newcommand{\ealkey}[1]{\textcolor{ealkeycolour}{#1}}
\newcommand{\ealkeytr}[1]{\textcolor{ealkeycolour}{\ealtext{#1}}}

% One translated word above one English word. Built from kernel
% primitives, never a tabular, which would break wherever the sheet
% loads array, booktabs or tabularx. The box is as wide as the wider of
% the two, so a long translation pushes its neighbours aside instead of
% overlapping them, and it claims its full height so a table row or a
% TikZ node makes room for it.
\newcommand{\ealgloss}[2]{%
  \sbox{\ealboxen}{\ealkey{#1}}%
  \sbox{\ealboxtr}{\small\textcolor{ealtwocolour}{\ealtext{#2}}}%
  \setlength{\ealglosswd}{\wd\ealboxen}%
  \ifdim\wd\ealboxtr>\ealglosswd\setlength{\ealglosswd}{\wd\ealboxtr}\fi
  \leavevmode
  \makebox[\ealglosswd][c]{%
    \makebox[0pt][c]{\usebox{\ealboxen}}%
    \makebox[0pt][c]{\raisebox{\ealglossraise}{\usebox{\ealboxtr}}}%
  }%
}

% opens up line spacing around a block of glosses, so the raised
% translations sit clear of the line above
\newenvironment{ealglossed}{\par\linespread{2.1}\selectfont\ignorespaces}{\par}

% a whole translated sentence above its English counterpart. block level,
% so it does not belong inside a tabular cell
\newcommand{\ealpara}[2]{%
  \par\addvspace{0.5\baselineskip}%
  {\color{ealtwocolour}\ealtextblock{#1}}%
  \penalty200
  {\color{ealenglish}#2\par}%
  \addvspace{0.5\baselineskip}%
}

% ================================================================
% Dictionary layout
% ================================================================
\usepackage{multicol}

\newlength{\ealdictgap}
\setlength{\ealdictgap}{1.2mm}

% english word / english meaning / translated word / translated meaning
\newcommand{\dictentrytr}[4]{%
  \par\addvspace{\ealdictgap}%
  \noindent
  {\bfseries\ealkey{#1}}\quad{\small #2}\par
  \nopagebreak
  {\color{ealtwocolour}\ealtext{{\bfseries #3}~\textemdash{} {\small #4}}\par}%
  \par
}

\pagestyle{empty}
\setlength{\parindent}{0pt}

\begin{document}
{\large\bfseries Punjabi Mathematical Dictionary}\par\medskip

\dictentrytr{absolute value}{the distance of a number from zero, ignoring whether it is positive or negative.}{ਨਿਰਪੇਖ ਮੁੱਲ}{ਕਿਸੇ ਸੰਖਿਆ ਦੀ ਸਿਫ਼ਰ ਤੋਂ ਦੂਰੀ, ਇਹ ਵੇਖੇ ਬਿਨਾਂ ਕਿ ਉਹ ਧਨਾਤਮਕ ਹੈ ਜਾਂ ਰਿਣਾਤਮਕ।}
\dictentrytr{addition}{combining two or more numbers to find their total.}{ਜੋੜ}{ਦੋ ਜਾਂ ਵੱਧ ਸੰਖਿਆਵਾਂ ਨੂੰ ਮਿਲਾ ਕੇ ਉਹਨਾਂ ਦਾ ਕੁੱਲ ਪਤਾ ਕਰਨਾ।}
\dictentrytr{adjacent}{next to, and sharing a side or a corner}{ਆਸਨਨ}{ਨਾਲ ਲੱਗਦਾ, ਅਤੇ ਇੱਕ ਭੁਜ ਜਾਂ ਕੋਨਾ ਸਾਂਝਾ ਕਰਦਾ ਹੋਇਆ}
\dictentrytr{algebraic method}{a way of solving a problem by using letters to stand for the numbers you want to find.}{ਬੀਜਗਣਿਤੀ ਵਿਧੀ}{ਕਿਸੇ ਸਮੱਸਿਆ ਨੂੰ ਹੱਲ ਕਰਨ ਦਾ ਇੱਕ ਤਰੀਕਾ ਜਿਸ ਵਿੱਚ ਅੱਖਰ ਉਹਨਾਂ ਸੰਖਿਆਵਾਂ ਨੂੰ ਦਰਸਾਉਣ ਲਈ ਵਰਤੇ ਜਾਂਦੇ ਹਨ ਜਿਨ੍ਹਾਂ ਨੂੰ ਤੁਸੀਂ ਲੱਭਣਾ ਚਾਹੁੰਦੇ ਹੋ।}
\dictentrytr{altitude}{height above sea level or another fixed reference level.}{ਉੱਚਾਈ}{ਸਮੁੰਦਰ ਤਲ ਜਾਂ ਕਿਸੇ ਹੋਰ ਸਥਿਰ ਸੰਦਰਭ ਪੱਧਰ ਤੋਂ ਉੱਪਰ ਦੀ ਉਚਾਈ।}
\dictentrytr{angle}{the amount of turn between two straight lines that meet at a point}{ਕੋਣ}{ਦੋ ਸਿੱਧੀਆਂ ਰੇਖਾਵਾਂ ਵਿਚਕਾਰ ਮੋੜ ਦੀ ਮਾਤਰਾ ਜੋ ਇੱਕ ਬਿੰਦੂ 'ਤੇ ਮਿਲਦੀਆਂ ਹਨ।}
\dictentrytr{annular}{Describing a shape made by the space between two concentric circles.}{ਵਲਯਾਕਾਰ}{ਇੱਕ ਸ਼ਕਲ ਦਾ ਵਰਣਨ ਕਰਨ ਵਾਲਾ ਜੋ ਦੋ ਸਮਕੇਂਦਰੀ ਚੱਕਰਾਂ ਵਿਚਕਾਰਲੀ ਥਾਂ ਨਾਲ ਬਣਦੀ ਹੈ।}
\dictentrytr{apex}{the point of a triangle opposite the base, often its top point.}{ਸਿਖਰ}{ਤਿਕੋਣ ਵਿੱਚ ਅਧਾਰ ਦੇ ਸਾਹਮਣੇ ਵਾਲਾ ਬਿੰਦੂ, ਅਕਸਰ ਉਸਦਾ ਸਭ ਤੋਂ ਉੱਪਰਲਾ ਬਿੰਦੂ।}
\dictentrytr{apex angle}{the angle at the point of a triangle opposite its base.}{ਸਿਖਰ ਕੋਣ}{ਤਿਕੋਣ ਦੇ ਆਧਾਰ ਦੇ ਸਾਹਮਣੇ ਵਾਲੇ ਸਿਖਰ 'ਤੇ ਬਣਿਆ ਕੋਣ।}
\dictentrytr{approximately}{close to but not exactly; about.}{ਲਗਭਗ}{ਨੇੜੇ ਪਰ ਬਿਲਕੁਲ ਨਹੀਂ; ਲਗਭਗ।}
\dictentrytr{arc}{part of the curved edge of a circle}{ਚਾਪ}{ਚੱਕਰ ਦੇ ਵਕਰੇ ਕਿਨਾਰੇ ਦਾ ਹਿੱਸਾ}
\dictentrytr{arccos}{The operation that finds an angle when its cosine is known.}{ਆਰਕਕੌਸ}{ਉਹ ਕਿਰਿਆ ਜੋ ਕਿਸੇ ਕੋਣ ਦਾ ਕੋਸਾਈਨ ਪਤਾ ਹੋਣ 'ਤੇ ਉਸ ਕੋਣ ਨੂੰ ਲੱਭਦੀ ਹੈ।}
\dictentrytr{arctan}{The operation that finds an angle when its tangent is known.}{ਆਰਕਟੈਨ}{ਉਹ ਕਿਰਿਆ ਜੋ ਕਿਸੇ ਕੋਣ ਦਾ ਟੈਨਜੈਂਟ ਪਤਾ ਹੋਣ 'ਤੇ ਉਸ ਕੋਣ ਨੂੰ ਲੱਭਦੀ ਹੈ।}
\dictentrytr{area}{the amount of space inside a flat shape}{ਖੇਤਰਫਲ}{ਕਿਸੇ ਸਮਤਲ ਆਕਾਰ ਦੇ ਅੰਦਰ ਥਾਂ ਦੀ ਮਾਤਰਾ}
\dictentrytr{area model}{a rectangle diagram split into smaller rectangles, used to show how to multiply fractions.}{ਖੇਤਰਫਲ ਮਾਡਲ}{ਇੱਕ ਆਇਤਾਕਾਰ ਚਿੱਤਰ ਜੋ ਛੋਟੇ ਆਇਤਾਂ ਵਿੱਚ ਵੰਡਿਆ ਜਾਂਦਾ ਹੈ, ਇਹ ਦਿਖਾਉਣ ਲਈ ਵਰਤਿਆ ਜਾਂਦਾ ਹੈ ਕਿ ਭਿੰਨਾਂ ਨੂੰ ਕਿਵੇਂ ਗੁਣਾ ਕਰਨਾ ਹੈ।}
\dictentrytr{arithmetic}{the branch of mathematics that deals with adding, subtracting, multiplying and dividing numbers.}{ਅੰਕਗਣਿਤ}{ਗਣਿਤ ਦੀ ਉਹ ਸ਼ਾਖਾ ਜੋ ਸੰਖਿਆਵਾਂ ਨੂੰ ਜੋੜਨ, ਘਟਾਉਣ, ਗੁਣਾ ਕਰਨ ਅਤੇ ਭਾਗ ਕਰਨ ਨਾਲ ਸੰਬੰਧਤ ਹੈ।}
\dictentrytr{arithmetic sequence}{a sequence where the same number is added or subtracted each time}{ਅੰਕਗਣਿਤਕ ਲੜੀ}{ਇੱਕ ਲੜੀ ਜਿਸ ਵਿੱਚ ਹਰ ਵਾਰ ਇੱਕੋ ਸੰਖਿਆ ਜੋੜੀ ਜਾਂ ਘਟਾਈ ਜਾਂਦੀ ਹੈ।}
\dictentrytr{ascending}{arranged from smallest to largest.}{ਚੜ੍ਹਦਾ ਕ੍ਰਮ}{ਛੋਟੇ ਤੋਂ ਵੱਡੇ ਕ੍ਰਮ ਵਿੱਚ ਵਿਵਸਥਿਤ।}
\dictentrytr{ascending order}{arranged from smallest to largest}{ਚੜ੍ਹਦਾ ਕ੍ਰਮ}{ਛੋਟੇ ਤੋਂ ਵੱਡੇ ਤੱਕ ਕ੍ਰਮਬੱਧ।}
\dictentrytr{average}{a single value chosen to stand for a whole set of data}{ਔਸਤ}{ਇੱਕ ਇੱਕਲਾ ਮੁੱਲ ਜੋ ਅੰਕੜਿਆਂ ਦੇ ਪੂਰੇ ਸਮੂਹ ਦੀ ਨੁਮਾਇੰਦਗੀ ਕਰਨ ਲਈ ਚੁਣਿਆ ਜਾਂਦਾ ਹੈ}
\dictentrytr{average speed}{total distance travelled divided by total time taken.}{ਔਸਤ ਗਤੀ}{ਤੈਅ ਕੀਤੀ ਕੁੱਲ ਦੂਰੀ ਨੂੰ ਲਏ ਗਏ ਕੁੱਲ ਸਮੇਂ ਨਾਲ ਭਾਗ ਕਰਨ ਦਾ ਨਤੀਜਾ।}
\dictentrytr{axis}{one of the lines a graph is measured against}{ਧੁਰਾ}{ਉਹਨਾਂ ਰੇਖਾਵਾਂ ਵਿੱਚੋਂ ਇੱਕ ਜਿਸ ਦੇ ਹਿਸਾਬ ਨਾਲ ਗ੍ਰਾਫ਼ ਨਾਪਿਆ ਜਾਂਦਾ ਹੈ}
\dictentrytr{bar model}{a diagram made of rectangles that shows the parts and total of an equation.}{ਬਾਰ ਮਾਡਲ}{ਇੱਕ ਚਿੱਤਰ ਜੋ ਆਇਤਾਂ ਨਾਲ ਬਣਦਾ ਹੈ ਅਤੇ ਕਿਸੇ ਸਮੀਕਰਨ ਦੇ ਭਾਗਾਂ ਅਤੇ ਕੁੱਲ ਨੂੰ ਦਰਸਾਉਂਦਾ ਹੈ।}
\dictentrytr{base}{the side of a shape chosen as the bottom for working out its area.}{ਆਧਾਰ}{ਕਿਸੇ ਸ਼ਕਲ ਦੀ ਉਹ ਭੁਜਾ ਜੋ ਉਸਦਾ ਖੇਤਰਫਲ ਕੱਢਣ ਲਈ ਥੱਲੇ ਵਜੋਂ ਚੁਣੀ ਜਾਂਦੀ ਹੈ।}
\dictentrytr{basis point}{a unit equal to one hundredth of one percent, used for small changes in rates.}{ਬੇਸਿਸ ਪੁਆਇੰਟ}{ਇੱਕ ਇਕਾਈ ਜੋ ਇੱਕ ਪ੍ਰਤੀਸ਼ਤ ਦੇ ਸੌਵੇਂ ਹਿੱਸੇ ਦੇ ਬਰਾਬਰ ਹੈ, ਦਰਾਂ ਵਿੱਚ ਛੋਟੀਆਂ ਤਬਦੀਲੀਆਂ ਲਈ ਵਰਤੀ ਜਾਂਦੀ ਹੈ।}
\dictentrytr{bearing}{a direction given as an angle clockwise from north, written with three figures}{ਬੇਅਰਿੰਗ}{ਇੱਕ ਦਿਸ਼ਾ ਜੋ ਉੱਤਰ ਤੋਂ ਘੜੀ ਦੀ ਦਿਸ਼ਾ ਵਿੱਚ ਕੋਣ ਦੇ ਰੂਪ ਵਿੱਚ ਦਿੱਤੀ ਜਾਂਦੀ ਹੈ, ਤਿੰਨ ਅੰਕਾਂ ਨਾਲ ਲਿਖੀ ਜਾਂਦੀ ਹੈ}
\dictentrytr{bimodal}{having two values that appear most often in a list of numbers.}{ਦੋ-ਬਹੁਲਕ}{ਅਜਿਹਾ ਜਿਸ ਵਿੱਚ ਸੰਖਿਆਵਾਂ ਦੀ ਸੂਚੀ ਵਿੱਚ ਦੋ ਮੁੱਲ ਸਭ ਤੋਂ ਵੱਧ ਵਾਰ ਆਉਂਦੇ ਹਨ।}
\dictentrytr{bisector}{a line that cuts an angle or another line exactly in half}{ਸਮਦੁਭਾਜਕ}{ਇੱਕ ਰੇਖਾ ਜੋ ਕਿਸੇ ਕੋਣ ਜਾਂ ਕਿਸੇ ਹੋਰ ਰੇਖਾ ਨੂੰ ਠੀਕ ਅੱਧਾ ਕੱਟਦੀ ਹੈ}
\dictentrytr{bodmas}{a rule that says which operation to do first in a calculation.}{ਬੋਡਮਾਸ}{ਇੱਕ ਨਿਯਮ ਜੋ ਦੱਸਦਾ ਹੈ ਕਿ ਕਿਸੇ ਗਣਨਾ ਵਿੱਚ ਕਿਹੜੀ ਕਿਰਿਆ ਪਹਿਲਾਂ ਕਰਨੀ ਹੈ।}
\dictentrytr{bond}{a type of loan to a government or company that pays regular money and is paid back later.}{ਬਾਂਡ}{ਇੱਕ ਕਿਸਮ ਦਾ ਕਰਜ਼ਾ ਜੋ ਸਰਕਾਰ ਜਾਂ ਕੰਪਨੀ ਨੂੰ ਦਿੱਤਾ ਜਾਂਦਾ ਹੈ, ਜੋ ਨਿਯਮਿਤ ਪੈਸੇ ਦਿੰਦਾ ਹੈ ਅਤੇ ਬਾਅਦ ਵਿੱਚ ਵਾਪਸ ਮੋੜਿਆ ਜਾਂਦਾ ਹੈ।}
\dictentrytr{box method}{a way of multiplying or expanding by splitting a rectangle into smaller parts.}{ਬਾਕਸ ਵਿਧੀ}{ਇੱਕ ਆਇਤ ਨੂੰ ਛੋਟੇ ਹਿੱਸਿਆਂ ਵਿੱਚ ਵੰਡ ਕੇ ਗੁਣਾ ਕਰਨ ਜਾਂ ਵਿਸਤਾਰ ਕਰਨ ਦਾ ਤਰੀਕਾ।}
\dictentrytr{bracket}{one of a pair of curved symbols, ( ), used to group numbers or letters.}{ਬਰੈਕਟ}{ਵਕਰ ਚਿੰਨ੍ਹਾਂ ( ) ਦੀ ਇੱਕ ਜੋੜੀ ਵਿੱਚੋਂ ਇੱਕ ਚਿੰਨ੍ਹ, ਜੋ ਸੰਖਿਆਵਾਂ ਜਾਂ ਅੱਖਰਾਂ ਨੂੰ ਸਮੂਹ ਕਰਨ ਲਈ ਵਰਤਿਆ ਜਾਂਦਾ ਹੈ।}
\dictentrytr{brackets}{A pair of signs used in mathematics to group together a part of an expression.}{ਬਰੈਕਟਾਂ}{ਗਣਿਤ ਵਿੱਚ ਵਰਤੇ ਜਾਣ ਵਾਲੇ ਚਿੰਨ੍ਹਾਂ ਦਾ ਇੱਕ ਜੋੜਾ ਜੋ ਕਿਸੇ ਵਿਅੰਜਕ ਦੇ ਹਿੱਸੇ ਨੂੰ ਇਕੱਠਾ ਰੱਖਣ ਲਈ ਵਰਤਿਆ ਜਾਂਦਾ ਹੈ।}
\dictentrytr{bus stop method}{a written method for division using a roof-like sign and a divisor.}{ਬੱਸ ਸਟਾਪ ਵਿਧੀ}{ਵੰਡ ਲਈ ਲਿਖਤੀ ਵਿਧੀ ਜਿਸ ਵਿੱਚ ਛੱਤ ਵਰਗਾ ਨਿਸ਼ਾਨ ਅਤੇ ਭਾਜਕ ਵਰਤਿਆ ਜਾਂਦਾ ਹੈ।}
\dictentrytr{cancel}{to remove a shared number from the top and bottom of a fraction before multiplying.}{ਰੱਦ ਕਰਨਾ}{ਗੁਣਾ ਕਰਨ ਤੋਂ ਪਹਿਲਾਂ ਕਿਸੇ ਭਿੰਨ ਦੇ ਉੱਪਰ ਅਤੇ ਹੇਠਾਂ ਤੋਂ ਸਾਂਝੀ ਸੰਖਿਆ ਨੂੰ ਹਟਾਉਣਾ।}
\dictentrytr{cardinal directions}{The four main compass points: North, South, East and West.}{ਮੁੱਖ ਦਿਸ਼ਾਵਾਂ}{ਕੰਪਾਸ ਦੇ ਚਾਰ ਮੁੱਖ ਬਿੰਦੂ: ਉੱਤਰ, ਦੱਖਣ, ਪੂਰਬ ਅਤੇ ਪੱਛਮ।}
\dictentrytr{category}{a group of non-numerical items, for example a colour, used as data.}{ਸ਼੍ਰੇਣੀ}{ਗੈਰ-ਸੰਖਿਅਕ ਵਸਤੂਆਂ ਦਾ ਸਮੂਹ, ਜਿਵੇਂ ਕਿ ਇੱਕ ਰੰਗ, ਜੋ ਅੰਕੜੇ ਵਜੋਂ ਵਰਤਿਆ ਜਾਂਦਾ ਹੈ।}
\dictentrytr{centre}{the fixed middle point of a circle}{ਕੇਂਦਰ}{ਚੱਕਰ ਦਾ ਸਥਿਰ ਮੱਧ ਬਿੰਦੂ।}
\dictentrytr{chord}{a straight line joining two points on the edge of a circle}{ਜੀਵਾ}{ਚੱਕਰ ਦੇ ਕਿਨਾਰੇ ਉੱਤੇ ਦੋ ਬਿੰਦੂਆਂ ਨੂੰ ਜੋੜਨ ਵਾਲੀ ਸਿੱਧੀ ਰੇਖਾ}
\dictentrytr{chunking}{a division method that breaks a larger number into easier smaller parts.}{ਚੰਕਿੰਗ}{ਇੱਕ ਭਾਗ ਕਰਨ ਦਾ ਤਰੀਕਾ ਜੋ ਵੱਡੀ ਸੰਖਿਆ ਨੂੰ ਛੋਟੇ ਅਸਾਨ ਹਿੱਸਿਆਂ ਵਿੱਚ ਤੋੜਦਾ ਹੈ।}
\dictentrytr{circle}{a round flat shape whose edge is always the same distance from its centre}{ਚੱਕਰ}{ਇੱਕ ਗੋਲ ਸਮਤਲ ਸ਼ਕਲ ਜਿਸਦਾ ਕੰਢਾ ਹਮੇਸ਼ਾ ਆਪਣੇ ਕੇਂਦਰ ਤੋਂ ਬਰਾਬਰ ਦੂਰੀ 'ਤੇ ਹੁੰਦਾ ਹੈ।}
\dictentrytr{circular}{round in shape, with every point on its edge the same distance from the centre.}{ਚੱਕਰਾਕਾਰ}{ਗੋਲ ਸ਼ਕਲ ਵਾਲਾ, ਜਿਸਦੇ ਕੰਢੇ ਦਾ ਹਰੇਕ ਬਿੰਦੂ ਕੇਂਦਰ ਤੋਂ ਬਰਾਬਰ ਦੂਰੀ 'ਤੇ ਹੁੰਦਾ ਹੈ।}
\dictentrytr{circumference}{the distance all the way round a circle}{ਪਰਿਧਿ}{ਚੱਕਰ ਦੇ ਚਾਰੇ ਪਾਸੇ ਦੀ ਦੂਰੀ}
\dictentrytr{coefficient}{the number in front of a letter, showing how many of it there are}{ਗੁਣਾਂਕ}{ਅੱਖਰ ਦੇ ਅੱਗੇ ਲੱਗੀ ਸੰਖਿਆ, ਜੋ ਦੱਸਦੀ ਹੈ ਕਿ ਉਸ ਦੀ ਕਿੰਨੀ ਮਾਤਰਾ ਹੈ।}
\dictentrytr{column vector}{a vector written vertically with one number above another.}{ਕਾਲਮ ਵੈਕਟਰ}{ਇੱਕ ਵੈਕਟਰ ਜੋ ਲੰਬਕਾਰੀ ਲਿਖਿਆ ਜਾਂਦਾ ਹੈ, ਜਿਸ ਵਿੱਚ ਇੱਕ ਸੰਖਿਆ ਦੂਜੀ ਦੇ ਉੱਪਰ ਹੁੰਦੀ ਹੈ।}
\dictentrytr{common denominator}{a bottom number shared by two or more fractions, used for adding or comparing them.}{ਸਾਂਝਾ ਹਰ}{ਦੋ ਜਾਂ ਵੱਧ ਭਿੰਨਾਂ ਵਿੱਚ ਸਾਂਝੀ ਹੇਠਲੀ ਸੰਖਿਆ, ਜੋ ਉਹਨਾਂ ਨੂੰ ਜੋੜਨ ਜਾਂ ਤੁਲਨਾ ਕਰਨ ਲਈ ਵਰਤੀ ਜਾਂਦੀ ਹੈ।}
\dictentrytr{common difference}{the fixed number added to each term to get the next term}{ਸਾਂਝਾ ਅੰਤਰ}{ਉਹ ਸਥਿਰ ਸੰਖਿਆ ਜੋ ਅਗਲਾ ਪਦ ਪ੍ਰਾਪਤ ਕਰਨ ਲਈ ਹਰ ਪਦ ਵਿੱਚ ਜੋੜੀ ਜਾਂਦੀ ਹੈ।}
\dictentrytr{common factor}{a number that divides exactly into two or more other numbers.}{ਸਾਂਝਾ ਗੁਣਨਖੰਡ}{ਉਹ ਸੰਖਿਆ ਜੋ ਦੋ ਜਾਂ ਵੱਧ ਹੋਰ ਸੰਖਿਆਵਾਂ ਨੂੰ ਪੂਰੀ ਤਰ੍ਹਾਂ ਵੰਡਦੀ ਹੈ।}
\dictentrytr{common ratio}{the fixed number used to create each next term in a geometric sequence}{ਸਾਂਝਾ ਅਨੁਪਾਤ}{ਉਹ ਸਥਿਰ ਸੰਖਿਆ ਜੋ ਗੁਣੋਤਰੀ ਲੜੀ ਵਿੱਚ ਹਰ ਅਗਲਾ ਪਦ ਬਣਾਉਣ ਲਈ ਵਰਤੀ ਜਾਂਦੀ ਹੈ।}
\dictentrytr{commutative}{describes an operation whose order does not affect the answer.}{ਕ੍ਰਮ-ਵਟਾਂਦਰਾ}{ਇੱਕ ਕਿਰਿਆ ਨੂੰ ਦਰਸਾਉਂਦਾ ਹੈ ਜਿਸਦਾ ਕ੍ਰਮ ਉੱਤਰ ਨੂੰ ਪ੍ਰਭਾਵਿਤ ਨਹੀਂ ਕਰਦਾ।}
\dictentrytr{commutativity}{the rule that swapping two numbers before you add or times them leaves the answer the same}{ਵਟਾਂਦਰਾ ਗੁਣ}{ਉਹ ਨਿਯਮ ਜੋ ਕਹਿੰਦਾ ਹੈ ਕਿ ਜੋੜ ਜਾਂ ਗੁਣਾ ਕਰਨ ਤੋਂ ਪਹਿਲਾਂ ਦੋ ਸੰਖਿਆਵਾਂ ਦੀ ਅਦਲ-ਬਦਲੀ ਕਰਨ ਨਾਲ ਜਵਾਬ ਇੱਕੋ ਜਿਹਾ ਰਹਿੰਦਾ ਹੈ।}
\dictentrytr{component}{one of the numbers that make up a vector.}{ਕੰਪੋਨੈਂਟ}{ਕਿਸੇ ਵੈਕਟਰ ਨੂੰ ਬਣਾਉਣ ਵਾਲੀਆਂ ਸੰਖਿਆਵਾਂ ਵਿੱਚੋਂ ਇੱਕ।}
\dictentrytr{compound}{to earn interest on both the original amount and on interest already earned.}{ਚੱਕਰਵਿੱਧੀ}{ਮੂਲ ਰਕਮ ਅਤੇ ਪਹਿਲਾਂ ਕਮਾਏ ਵਿਆਜ ਦੋਵਾਂ 'ਤੇ ਵਿਆਜ ਕਮਾਉਣਾ।}
\dictentrytr{compounding}{the process of earning profit on both the original amount and profit already earned.}{ਚੱਕਰਵਰਧੀ}{ਮੂਲ ਰਕਮ ਅਤੇ ਪਹਿਲਾਂ ਕਮਾਏ ਲਾਭ ਦੋਵਾਂ 'ਤੇ ਲਾਭ ਕਮਾਉਣ ਦੀ ਪ੍ਰਕਿਰਿਆ।}
\dictentrytr{concentric}{Having the same centre.}{ਸਮਕੇਂਦਰੀ}{ਇੱਕੋ ਕੇਂਦਰ ਵਾਲਾ।}
\dictentrytr{congruent}{exactly the same shape and size}{ਸਰਬੰਗਸਮ}{ਬਿਲਕੁਲ ਸਮਾਨ ਸ਼ਕਲ ਅਤੇ ਆਕਾਰ ਵਾਲਾ}
\dictentrytr{constant}{a fixed number that stands alone, like 3 in 3x+3.}{ਸਥਿਰਾਂਕ}{ਇੱਕ ਸਥਿਰ ਸੰਖਿਆ ਜੋ ਇਕੱਲੀ ਖੜ੍ਹੀ ਹੁੰਦੀ ਹੈ, ਜਿਵੇਂ 3x+3 ਵਿੱਚ 3।}
\dictentrytr{convergence}{the process of getting closer and closer to a particular value}{ਅਭਿਸਰਣ}{ਕਿਸੇ ਖਾਸ ਮੁੱਲ ਦੇ ਹੋਰ ਅਤੇ ਹੋਰ ਨੇੜੇ ਹੁੰਦੇ ਜਾਣ ਦੀ ਪ੍ਰਕਿਰਿਆ।}
\dictentrytr{conversion}{a change of a value from one form, such as a fraction, to another, such as a percentage.}{ਰੂਪਾਂਤਰਨ}{ਕਿਸੇ ਮੁੱਲ ਦਾ ਇੱਕ ਰੂਪ ਤੋਂ ਦੂਜੇ ਰੂਪ ਵਿੱਚ ਬਦਲਣਾ, ਜਿਵੇਂ ਕਿ ਭਿੰਨ ਤੋਂ ਪ੍ਰਤੀਸ਼ਤ ਵਿੱਚ।}
\dictentrytr{convert}{to change a number from one form into another, such as a mixed number into an improper fraction}{ਰੂਪਾਂਤਰਿਤ ਕਰਨਾ}{ਕਿਸੇ ਸੰਖਿਆ ਨੂੰ ਇੱਕ ਰੂਪ ਤੋਂ ਦੂਜੇ ਰੂਪ ਵਿੱਚ ਬਦਲਣਾ, ਜਿਵੇਂ ਕਿ ਮਿਸ਼ਰਿਤ ਸੰਖਿਆ ਨੂੰ ਅਨੁਚਿਤ ਭਿੰਨ ਵਿੱਚ।}
\dictentrytr{coordinate}{a number or pair of numbers that fixes the position of a point.}{ਨਿਰਦੇਸ਼ਾਂਕ}{ਇੱਕ ਸੰਖਿਆ ਜਾਂ ਸੰਖਿਆਵਾਂ ਦਾ ਜੋੜਾ ਜੋ ਕਿਸੇ ਬਿੰਦੂ ਦੀ ਸਥਿਤੀ ਨੂੰ ਨਿਸ਼ਚਿਤ ਕਰਦਾ ਹੈ।}
\dictentrytr{correlation}{how closely two sets of data are related}{ਸਹ-ਸਬੰਧ}{ਅੰਕੜਿਆਂ ਦੇ ਦੋ ਸਮੂਹ ਕਿੰਨੇ ਨੇੜਿਓਂ ਸੰਬੰਧਿਤ ਹਨ}
\dictentrytr{corresponding}{matching sides or points in similar shapes.}{ਸੰਗਤ}{ਸਮਰੂਪ ਆਕਾਰਾਂ ਵਿੱਚ ਮਿਲਦੀਆਂ ਭੁਜਾਵਾਂ ਜਾਂ ਬਿੰਦੂ।}
\dictentrytr{cos}{the adjacent side divided by the hypotenuse for an angle in a right triangle.}{ਕੋਸ}{ਸਮਕੋਣ ਤਿਕੋਣ ਵਿੱਚ ਕਿਸੇ ਕੋਣ ਲਈ ਲਾਗਲੀ ਭੁਜਾ ਨੂੰ ਕਰਣ ਨਾਲ ਵੰਡਿਆ ਜਾਂਦਾ ਹੈ।}
\dictentrytr{cosine}{in a three-sided shape with a right angle, the ratio of the adjacent side to the hypotenuse.}{ਕੋਸਾਈਨ}{ਸਮਕੋਣ ਤਿਕੋਣ ਵਿੱਚ, ਲਾਗਲੀ ਭੁਜਾ ਅਤੇ ਕਰਣ ਦਾ ਅਨੁਪਾਤ।}
\dictentrytr{cosine rule}{a formula connecting three side lengths and one angle in a three-sided shape.}{ਕੋਸਾਈਨ ਨਿਯਮ}{ਤਿਕੋਣ ਵਿੱਚ ਤਿੰਨ ਭੁਜਾਵਾਂ ਦੀਆਂ ਲੰਬਾਈਆਂ ਅਤੇ ਇੱਕ ਕੋਣ ਨੂੰ ਜੋੜਨ ਵਾਲਾ ਸੂਤਰ।}
\dictentrytr{counterexample}{an example that proves a general statement is not always true}{ਪ੍ਰਤਿਉਦਾਹਰਣ}{ਇੱਕ ਉਦਾਹਰਣ ਜੋ ਦਰਸਾਉਂਦੀ ਹੈ ਕਿ ਇੱਕ ਆਮ ਬਿਆਨ ਹਮੇਸ਼ਾ ਸੱਚ ਨਹੀਂ ਹੁੰਦਾ।}
\dictentrytr{coupon}{The regular interest payment made on a bond.}{ਕੂਪਨ}{ਬਾਂਡ ਉੱਤੇ ਕੀਤਾ ਜਾਣ ਵਾਲਾ ਨਿਯਮਤ ਵਿਆਜ ਭੁਗਤਾਨ।}
\dictentrytr{cube}{the number you get when a value is multiplied by itself twice.}{ਘਣ}{ਉਹ ਸੰਖਿਆ ਜੋ ਕਿਸੇ ਮੁੱਲ ਨੂੰ ਆਪਣੇ ਨਾਲ ਦੋ ਵਾਰ ਗੁਣਾ ਕਰਨ ਨਾਲ ਪ੍ਰਾਪਤ ਹੁੰਦੀ ਹੈ।}
\dictentrytr{cuboid}{a 3D solid shape whose six faces are all rectangles.}{ਘਣਾਭ}{ਇੱਕ ਤਿੰਨ-ਅਯਾਮੀ ਠੋਸ ਆਕਾਰ ਜਿਸ ਦੇ ਛੇ ਫਲਕ ਸਾਰੇ ਆਇਤ ਹੁੰਦੇ ਹਨ।}
\dictentrytr{cumulative}{adding up as you go along, so each total includes the ones before}{ਸੰਚਿਤ}{ਅੱਗੇ ਵਧਦੇ ਹੋਏ ਜੋੜਦੇ ਜਾਣਾ, ਤਾਂ ਕਿ ਹਰ ਕੁੱਲ ਜੋੜ ਵਿੱਚ ਪਿਛਲੇ ਸ਼ਾਮਲ ਹੋਣ}
\dictentrytr{data}{the set of values collected}{ਅੰਕੜੇ}{ਇਕੱਠੇ ਕੀਤੇ ਗਏ ਮੁੱਲਾਂ ਦਾ ਸਮੂਹ}
\dictentrytr{dataset}{a collection of numbers or information being studied}{ਡਾਟਾਸੈੱਟ}{ਅਧਿਐਨ ਕੀਤੇ ਜਾ ਰਹੇ ਅੰਕੜਿਆਂ ਜਾਂ ਜਾਣਕਾਰੀ ਦਾ ਸੰਗ੍ਰਹਿ।}
\dictentrytr{decimal}{a number written with a point, where the digits after the point are parts of a whole}{ਦਸ਼ਮਲਵ}{ਇੱਕ ਸੰਖਿਆ ਜੋ ਬਿੰਦੂ ਨਾਲ ਲਿਖੀ ਜਾਂਦੀ ਹੈ, ਜਿੱਥੇ ਬਿੰਦੂ ਤੋਂ ਬਾਅਦ ਦੇ ਅੰਕ ਪੂਰੇ ਦੇ ਭਾਗ ਹੁੰਦੇ ਹਨ।}
\dictentrytr{decimal place}{the position of a digit after the decimal point, such as tenths or hundredths}{ਦਸ਼ਮਲਵ ਸਥਾਨ}{ਦਸ਼ਮਲਵ ਬਿੰਦੂ ਤੋਂ ਬਾਅਦ ਕਿਸੇ ਅੰਕ ਦੀ ਸਥਿਤੀ, ਜਿਵੇਂ ਦਸਵਾਂ ਜਾਂ ਸੌਵਾਂ।}
\dictentrytr{decimal places}{the digits written after the point in a number.}{ਦਸ਼ਮਲਵ ਸਥਾਨ}{ਕਿਸੇ ਸੰਖਿਆ ਵਿੱਚ ਦਸ਼ਮਲਵ ਬਿੰਦੂ ਤੋਂ ਬਾਅਦ ਲਿਖੇ ਅੰਕ।}
\dictentrytr{decimal point}{the dot used to separate the whole part of a number from the fractional part.}{ਦਸ਼ਮਲਵ ਬਿੰਦੂ}{ਉਹ ਬਿੰਦੀ ਜੋ ਕਿਸੇ ਸੰਖਿਆ ਦੇ ਪੂਰੇ ਹਿੱਸੇ ਨੂੰ ਭਿੰਨਾਤਮਕ ਹਿੱਸੇ ਤੋਂ ਵੱਖ ਕਰਨ ਲਈ ਵਰਤੀ ਜਾਂਦੀ ਹੈ।}
\dictentrytr{degree}{a unit for measuring angles; 360 of them make one full turn.}{ਡਿਗਰੀ}{ਕੋਣਾਂ ਨੂੰ ਮਾਪਣ ਦੀ ਇਕਾਈ; ਇਨ੍ਹਾਂ ਵਿੱਚੋਂ 360 ਇੱਕ ਪੂਰਾ ਚੱਕਰ ਬਣਾਉਂਦੇ ਹਨ।}
\dictentrytr{denominator}{the number below the line in a fraction, showing how many equal parts the whole is split into}{ਹਰ}{ਭਿੰਨ ਵਿੱਚ ਰੇਖਾ ਦੇ ਹੇਠਾਂ ਵਾਲੀ ਸੰਖਿਆ, ਜੋ ਦੱਸਦੀ ਹੈ ਕਿ ਪੂਰਾ ਕਿੰਨੇ ਬਰਾਬਰ ਹਿੱਸਿਆਂ ਵਿੱਚ ਵੰਡਿਆ ਗਿਆ ਹੈ।}
\dictentrytr{depreciation}{the loss in value of something over time.}{ਮੁੱਲ ਘਾਟਾ}{ਸਮੇਂ ਦੇ ਨਾਲ ਕਿਸੇ ਚੀਜ਼ ਦੇ ਮੁੱਲ ਵਿੱਚ ਹੋਣ ਵਾਲਾ ਨੁਕਸਾਨ।}
\dictentrytr{diameter}{a straight line across a circle through its centre}{ਵਿਆਸ}{ਚੱਕਰ ਦੇ ਕੇਂਦਰ ਵਿੱਚੋਂ ਲੰਘਦੀ ਹੋਈ ਇੱਕ ਪਾਸਿਓਂ ਦੂਜੇ ਪਾਸੇ ਤੱਕ ਸਿੱਧੀ ਰੇਖਾ}
\dictentrytr{difference}{the answer when one number is subtracted from another}{ਅੰਤਰ}{ਇੱਕ ਸੰਖਿਆ ਵਿੱਚੋਂ ਦੂਜੀ ਨੂੰ ਘਟਾਉਣ 'ਤੇ ਮਿਲਣ ਵਾਲਾ ਉੱਤਰ।}
\dictentrytr{digit}{a single figure from 0 to 9 that a number is written with}{ਅੰਕ}{ਇੱਕ ਅੰਕ 0 ਤੋਂ 9 ਤੱਕ, ਜੋ ਕਿਸੇ ਸੰਖਿਆ ਨੂੰ ਲਿਖਣ ਲਈ ਵਰਤਿਆ ਜਾਂਦਾ ਹੈ।}
\dictentrytr{dimension}{a measurement of a shape, such as its length or width.}{ਅਯਾਮ}{ਕਿਸੇ ਆਕਾਰ ਦਾ ਮਾਪ, ਜਿਵੇਂ ਉਸਦੀ ਲੰਬਾਈ ਜਾਂ ਚੌੜਾਈ।}
\dictentrytr{dimensions}{the measurements of a shape, such as length, width and height.}{ਅਯਾਮ}{ਕਿਸੇ ਆਕਾਰ ਦੇ ਮਾਪ, ਜਿਵੇਂ ਲੰਬਾਈ, ਚੌੜਾਈ ਅਤੇ ਉੱਚਾਈ।}
\dictentrytr{direct proportion}{a relationship where two amounts change together while their ratio stays the same.}{ਸਿੱਧਾ ਅਨੁਪਾਤ}{ਇੱਕ ਅਜਿਹਾ ਸੰਬੰਧ ਜਿੱਥੇ ਦੋ ਮਾਤਰਾਵਾਂ ਇਕੱਠੀਆਂ ਬਦਲਦੀਆਂ ਹਨ ਅਤੇ ਉਨ੍ਹਾਂ ਦਾ ਅਨੁਪਾਤ ਇੱਕੋ ਜਿਹਾ ਰਹਿੰਦਾ ਹੈ।}
\dictentrytr{directed path}{a line with an arrow showing direction from one place to another}{ਨਿਰਦੇਸ਼ਿਤ ਮਾਰਗ}{ਇੱਕ ਲਕੀਰ ਜਿਸ ਉੱਤੇ ਤੀਰ ਇੱਕ ਸਥਾਨ ਤੋਂ ਦੂਜੇ ਸਥਾਨ ਦੀ ਦਿਸ਼ਾ ਦਰਸਾਉਂਦਾ ਹੈ।}
\dictentrytr{directly proportional}{when one amount is multiplied by a number, the other amount is multiplied by that same number.}{ਸਿੱਧਾ ਅਨੁਪਾਤੀ}{ਜਦੋਂ ਇੱਕ ਮਾਤਰਾ ਨੂੰ ਕਿਸੇ ਸੰਖਿਆ ਨਾਲ ਗੁਣਾ ਕੀਤਾ ਜਾਂਦਾ ਹੈ, ਤਾਂ ਦੂਜੀ ਮਾਤਰਾ ਨੂੰ ਉਸੇ ਸੰਖਿਆ ਨਾਲ ਗੁਣਾ ਕੀਤਾ ਜਾਂਦਾ ਹੈ।}
\dictentrytr{discriminant}{a value calculated from the numbers in an equation; if it is negative, the equation has no real solutions.}{ਵਿਵੇਚਕ}{ਕਿਸੇ ਸਮੀਕਰਨ ਵਿੱਚ ਸੰਖਿਆਵਾਂ ਤੋਂ ਗਿਣਿਆ ਗਿਆ ਮੁੱਲ; ਜੇ ਇਹ ਨਕਾਰਾਤਮਕ ਹੋਵੇ, ਤਾਂ ਸਮੀਕਰਨ ਦਾ ਕੋਈ ਅਸਲ ਹੱਲ ਨਹੀਂ ਹੁੰਦਾ।}
\dictentrytr{displacement}{the straight-line distance and direction from a starting point to a final position.}{ਵਿਸਥਾਪਨ}{ਕਿਸੇ ਸ਼ੁਰੂਆਤੀ ਬਿੰਦੂ ਤੋਂ ਅੰਤਿਮ ਸਥਿਤੀ ਤੱਕ ਦੀ ਸਿੱਧੀ-ਲਕੀਰ ਦੂਰੀ ਅਤੇ ਦਿਸ਼ਾ।}
\dictentrytr{distributive property}{the rule that a(b+c)=ab+ac.}{ਵਿਤਰਨ ਗੁਣ}{ਉਹ ਨਿਯਮ ਜੋ ਦੱਸਦਾ ਹੈ ਕਿ a(b+c)=ab+ac.}
\dictentrytr{divide}{to split a number into equal parts, or find how many times one number fits inside another.}{ਵੰਡਣਾ}{ਕਿਸੇ ਸੰਖਿਆ ਨੂੰ ਬਰਾਬਰ ਹਿੱਸਿਆਂ ਵਿੱਚ ਵੰਡਣਾ, ਜਾਂ ਇਹ ਪਤਾ ਕਰਨਾ ਕਿ ਇੱਕ ਸੰਖਿਆ ਦੂਜੀ ਵਿੱਚ ਕਿੰਨੀ ਵਾਰ ਸਮਾਉਂਦੀ ਹੈ।}
\dictentrytr{division}{the operation of splitting a number into equal parts.}{ਭਾਗ}{ਕਿਸੇ ਸੰਖਿਆ ਨੂੰ ਬਰਾਬਰ ਹਿੱਸਿਆਂ ਵਿੱਚ ਵੰਡਣ ਦੀ ਕਿਰਿਆ।}
\dictentrytr{divisor}{the number that another number is divided by.}{ਭਾਜਕ}{ਉਹ ਸੰਖਿਆ ਜਿਸ ਨਾਲ ਕਿਸੇ ਹੋਰ ਸੰਖਿਆ ਨੂੰ ਵੰਡਿਆ ਜਾਂਦਾ ਹੈ।}
\dictentrytr{duration}{the length of time until a bond or gilt is paid back.}{ਮਿਆਦ}{ਬਾਂਡ ਜਾਂ ਗਿਲਟ ਦੇ ਵਾਪਸ ਭੁਗਤਾਨ ਹੋਣ ਤੱਕ ਦਾ ਸਮਾਂ।}
\dictentrytr{edge}{a straight line where two faces of a solid meet}{ਕਿਨਾਰਾ}{ਇੱਕ ਸਿੱਧੀ ਰੇਖਾ ਜਿੱਥੇ ਕਿਸੇ ਠੋਸ ਆਕਾਰ ਦੀਆਂ ਦੋ ਫਲਕਾਂ ਮਿਲਦੀਆਂ ਹਨ}
\dictentrytr{elimination}{removing one unknown from simultaneous equations by adding or subtracting them}{ਵਿਲੋਪਨ}{ਸਮਕਾਲੀ ਸਮੀਕਰਨਾਂ ਵਿੱਚੋਂ ਇੱਕ ਅਣਜਾਣ ਰਾਸ਼ੀ ਨੂੰ ਉਹਨਾਂ ਨੂੰ ਜੋੜ ਕੇ ਜਾਂ ਘਟਾ ਕੇ ਹਟਾਉਣਾ।}
\dictentrytr{embedding}{a way of turning words into coordinates or vectors so meanings can be compared mathematically}{ਏਮਬੈਡਿੰਗ}{ਸ਼ਬਦਾਂ ਨੂੰ ਨਿਰਦੇਸ਼ਾਂਕਾਂ ਜਾਂ ਵੈਕਟਰਾਂ ਵਿੱਚ ਬਦਲਣ ਦਾ ਤਰੀਕਾ ਤਾਂ ਜੋ ਅਰਥਾਂ ਦੀ ਗਣਿਤਕ ਤੌਰ 'ਤੇ ਤੁਲਨਾ ਕੀਤੀ ਜਾ ਸਕੇ।}
\dictentrytr{endpoint}{the point at which a route or line finishes}{ਅੰਤ ਬਿੰਦੂ}{ਉਹ ਬਿੰਦੂ ਜਿੱਥੇ ਕੋਈ ਰਸਤਾ ਜਾਂ ਰੇਖਾ ਖਤਮ ਹੁੰਦੀ ਹੈ।}
\dictentrytr{enlarge}{to make a shape bigger by multiplying every length by the same number.}{ਵੱਡਾ ਕਰਨਾ}{ਹਰ ਲੰਬਾਈ ਨੂੰ ਇੱਕੋ ਸੰਖਿਆ ਨਾਲ ਗੁਣਾ ਕਰਕੇ ਕਿਸੇ ਆਕਾਰ ਨੂੰ ਵੱਡਾ ਕਰਨਾ।}
\dictentrytr{equal}{having the same size or number}{ਬਰਾਬਰ}{ਇੱਕੋ ਆਕਾਰ ਜਾਂ ਇੱਕੋ ਸੰਖਿਆ ਵਾਲਾ।}
\dictentrytr{equation}{a mathematical statement with an equals sign showing that two sides are equal.}{ਸਮੀਕਰਨ}{ਇੱਕ ਗਣਿਤਕ ਬਿਆਨ ਜਿਸ ਵਿੱਚ ਬਰਾਬਰ ਦਾ ਨਿਸ਼ਾਨ ਹੁੰਦਾ ਹੈ ਜੋ ਦਰਸਾਉਂਦਾ ਹੈ ਕਿ ਦੋਵੇਂ ਪਾਸੇ ਬਰਾਬਰ ਹਨ।}
\dictentrytr{equilateral}{having all sides the same length}{ਸਮਭੁਜ}{ਜਿਸ ਦੀਆਂ ਸਾਰੀਆਂ ਭੁਜਾਵਾਂ ਦੀ ਲੰਬਾਈ ਇੱਕੋ ਜਿਹੀ ਹੋਵੇ।}
\dictentrytr{equilateral triangle}{a triangle with all three sides the same length.}{ਸਮਭੁਜ ਤਿਕੋਣ}{ਇੱਕ ਤਿਕੋਣ ਜਿਸ ਦੀਆਂ ਤਿੰਨੋਂ ਭੁਜਾਵਾਂ ਦੀ ਲੰਬਾਈ ਇੱਕੋ ਜਿਹੀ ਹੋਵੇ।}
\dictentrytr{equivalent}{having the same value, even though it is written differently}{ਤੁੱਲ}{ਮੁੱਲ ਇੱਕੋ ਜਿਹਾ ਹੋਣਾ, ਭਾਵੇਂ ਵੱਖਰੇ ਰੂਪ ਵਿੱਚ ਲਿਖਿਆ ਗਿਆ ਹੋਵੇ।}
\dictentrytr{equivalent fraction}{a different fraction that means the same amount as another fraction.}{ਤੁੱਲ ਭਿੰਨ}{ਇੱਕ ਵੱਖਰੀ ਭਿੰਨ ਜੋ ਕਿਸੇ ਹੋਰ ਭਿੰਨ ਦੇ ਬਰਾਬਰ ਮਾਤਰਾ ਨੂੰ ਦਰਸਾਉਂਦੀ ਹੈ।}
\dictentrytr{estimate}{a sensible answer worked out quickly from rounded numbers}{ਅੰਦਾਜ਼ਾ}{ਗੋਲ ਕੀਤੀਆਂ ਸੰਖਿਆਵਾਂ ਤੋਂ ਜਲਦੀ ਕੱਢਿਆ ਗਿਆ ਵਾਜਬ ਉੱਤਰ।}
\dictentrytr{evaluate}{to work out the value of an expression}{ਮੁੱਲ ਕੱਢਣਾ}{ਕਿਸੇ ਗਣਿਤਕ ਰਾਸ਼ੀ ਦਾ ਮੁੱਲ ਪਤਾ ਕਰਨਾ।}
\dictentrytr{even}{able to be divided exactly by 2.}{ਸਮ}{ਜੋ 2 ਨਾਲ ਪੂਰੀ ਤਰ੍ਹਾਂ ਵੰਡਿਆ ਜਾ ਸਕੇ।}
\dictentrytr{even number}{a whole number that can be divided exactly by two.}{ਸਮ ਸੰਖਿਆ}{ਉਹ ਪੂਰਨ ਸੰਖਿਆ ਜਿਸਨੂੰ ਦੋ ਨਾਲ ਪੂਰੀ ਤਰ੍ਹਾਂ ਵੰਡਿਆ ਜਾ ਸਕੇ।}
\dictentrytr{expand}{to multiply out brackets so the expression has no brackets left}{ਖੋਲ੍ਹਣਾ}{ਬਰੈਕਟਾਂ ਦਾ ਗੁਣਾ ਕਰਕੇ ਉਨ੍ਹਾਂ ਨੂੰ ਖੋਲ੍ਹਣਾ ਤਾਂ ਜੋ ਵਿਅੰਜਕ ਵਿੱਚ ਕੋਈ ਬਰੈਕਟ ਨਾ ਬਚੇ।}
\dictentrytr{exponent}{the raised number showing how many times to multiply a number by itself.}{ਘਾਤ}{ਉੱਪਰ ਲਿਖੀ ਸੰਖਿਆ ਜੋ ਦਰਸਾਉਂਦੀ ਹੈ ਕਿ ਕਿਸੇ ਸੰਖਿਆ ਨੂੰ ਆਪਣੇ ਨਾਲ ਕਿੰਨੀ ਵਾਰ ਗੁਣਾ ਕਰਨਾ ਹੈ।}
\dictentrytr{expression}{a group of numbers and letters joined by operations, with no equals sign}{ਵਿਅੰਜਕ}{ਸੰਖਿਆਵਾਂ ਅਤੇ ਅੱਖਰਾਂ ਦਾ ਸਮੂਹ ਜੋ ਓਪਰੇਸ਼ਨਾਂ ਨਾਲ ਜੁੜਿਆ ਹੋਵੇ, ਜਿਸ ਵਿੱਚ ਬਰਾਬਰ ਦਾ ਨਿਸ਼ਾਨ ਨਾ ਹੋਵੇ।}
\dictentrytr{face}{a flat surface of a solid shape}{ਫਲਕ}{ਕਿਸੇ ਠੋਸ ਆਕਾਰ ਦੀ ਇੱਕ ਸਮਤਲ ਸਤਹ}
\dictentrytr{factor}{a whole number that divides exactly into another number}{ਗੁਣਨਖੰਡ}{ਇੱਕ ਪੂਰਨ ਅੰਕ ਜੋ ਕਿਸੇ ਹੋਰ ਸੰਖਿਆ ਨੂੰ ਪੂਰਾ-ਪੂਰਾ ਵੰਡਦਾ ਹੈ।}
\dictentrytr{factorise}{to write an expression as a product, usually by putting brackets back in}{ਗੁਣਨਖੰਡ ਕਰਨਾ}{ਕਿਸੇ ਵਿਅੰਜਕ ਨੂੰ ਗੁਣਨਫਲ ਦੇ ਰੂਪ ਵਿੱਚ ਲਿਖਣਾ, ਆਮ ਤੌਰ 'ਤੇ ਬਰੈਕਟਾਂ ਵਾਪਸ ਲਗਾ ਕੇ।}
\dictentrytr{fibonacci}{a famous sequence where each number is the total of the two before it}{ਫਿਬੋਨਾਚੀ ਲੜੀ}{ਇੱਕ ਮਸ਼ਹੂਰ ਲੜੀ ਜਿਸ ਵਿੱਚ ਹਰ ਸੰਖਿਆ ਆਪਣੇ ਤੋਂ ਪਹਿਲਾਂ ਦੀਆਂ ਦੋ ਸੰਖਿਆਵਾਂ ਦਾ ਜੋੜ ਹੁੰਦੀ ਹੈ।}
\dictentrytr{finite}{limited in size; not going on forever.}{ਸੀਮਤ}{ਆਕਾਰ ਵਿੱਚ ਸੀਮਤ; ਹਮੇਸ਼ਾ ਜਾਰੀ ਨਾ ਰਹਿਣ ਵਾਲਾ।}
\dictentrytr{fixed point}{a value that stays the same when an iterative process is applied to it.}{ਸਥਿਰ ਬਿੰਦੂ}{ਇੱਕ ਮੁੱਲ ਜੋ ਉਦੋਂ ਇੱਕੋ ਜਿਹਾ ਰਹਿੰਦਾ ਹੈ ਜਦੋਂ ਉਸ ਉੱਤੇ ਪੁਨਰਾਵਰਤੀ ਪ੍ਰਕਿਰਿਆ ਲਾਗੂ ਕੀਤੀ ਜਾਂਦੀ ਹੈ।}
\dictentrytr{foil}{a way to multiply two brackets that each have two terms, multiplying First, Outer, Inner and Last.}{ਫੋਇਲ}{ਦੋ ਬਰੈਕਟਾਂ ਨੂੰ ਗੁਣਾ ਕਰਨ ਦਾ ਤਰੀਕਾ, ਜਿਸ ਵਿੱਚ ਹਰੇਕ ਦੇ ਦੋ ਪਦ ਹੋਣ, ਪਹਿਲੇ, ਬਾਹਰਲੇ, ਅੰਦਰਲੇ ਅਤੇ ਆਖਰੀ ਪਦਾਂ ਨੂੰ ਗੁਣਾ ਕੀਤਾ ਜਾਂਦਾ ਹੈ।}
\dictentrytr{formula}{a rule written with letters, showing how to work one quantity out from others}{ਸੂਤਰ}{ਅੱਖਰਾਂ ਨਾਲ ਲਿਖਿਆ ਨਿਯਮ, ਜੋ ਦੱਸਦਾ ਹੈ ਕਿ ਇੱਕ ਮਾਤਰਾ ਨੂੰ ਦੂਜੀਆਂ ਤੋਂ ਕਿਵੇਂ ਕੱਢਣਾ ਹੈ।}
\dictentrytr{fraction}{a number written as one whole split into equal parts, with a number of those parts taken}{ਭਿੰਨ}{ਇੱਕ ਸੰਖਿਆ ਜੋ ਇੱਕ ਪੂਰੇ ਨੂੰ ਬਰਾਬਰ ਹਿੱਸਿਆਂ ਵਿੱਚ ਵੰਡ ਕੇ, ਉਨ੍ਹਾਂ ਵਿੱਚੋਂ ਕੁਝ ਹਿੱਸੇ ਲੈ ਕੇ ਲਿਖੀ ਜਾਂਦੀ ਹੈ।}
\dictentrytr{fraction wall}{a diagram of strips divided into equal parts, used to compare fractions.}{ਭਿੰਨ ਕੰਧ}{ਪੱਟੀਆਂ ਦਾ ਚਿੱਤਰ ਜੋ ਬਰਾਬਰ ਹਿੱਸਿਆਂ ਵਿੱਚ ਵੰਡਿਆ ਹੁੰਦਾ ਹੈ, ਭਿੰਨਾਂ ਦੀ ਤੁਲਨਾ ਕਰਨ ਲਈ ਵਰਤਿਆ ਜਾਂਦਾ ਹੈ।}
\dictentrytr{fractional part}{the part of a mixed number that is less than one, written as a fraction.}{ਭਿੰਨਾਤਮਕ ਭਾਗ}{ਮਿਸ਼ਰਤ ਸੰਖਿਆ ਦਾ ਉਹ ਭਾਗ ਜੋ ਇੱਕ ਤੋਂ ਘੱਟ ਹੁੰਦਾ ਹੈ, ਭਿੰਨ ਵਜੋਂ ਲਿਖਿਆ ਜਾਂਦਾ ਹੈ।}
\dictentrytr{frequency}{how many times something happens}{ਬਾਰੰਬਾਰਤਾ}{ਕੋਈ ਚੀਜ਼ ਕਿੰਨੀ ਵਾਰ ਵਾਪਰਦੀ ਹੈ}
\dictentrytr{geometric}{describing a sequence where each term is multiplied by the same number to get the next.}{ਗੁਣੋਤਰੀ}{ਇੱਕ ਲੜੀ ਦਾ ਵਰਣਨ ਜਿੱਥੇ ਹਰੇਕ ਪਦ ਨੂੰ ਅਗਲਾ ਪਦ ਪ੍ਰਾਪਤ ਕਰਨ ਲਈ ਉਸੇ ਸੰਖਿਆ ਨਾਲ ਗੁਣਾ ਕੀਤਾ ਜਾਂਦਾ ਹੈ।}
\dictentrytr{geometric sequence}{a sequence where each term is found by multiplying the previous term by the same number}{ਗੁਣੋਤਰੀ ਲੜੀ}{ਇੱਕ ਲੜੀ ਜਿਸ ਵਿੱਚ ਹਰ ਪਦ ਪਿਛਲੇ ਪਦ ਨੂੰ ਇੱਕੋ ਸੰਖਿਆ ਨਾਲ ਗੁਣਾ ਕਰਕੇ ਪ੍ਰਾਪਤ ਕੀਤਾ ਜਾਂਦਾ ਹੈ।}
\dictentrytr{geometry}{the branch of mathematics that studies shapes, sizes, and space.}{ਰੇਖਾਗਣਿਤ}{ਗਣਿਤ ਦੀ ਸ਼ਾਖਾ ਜੋ ਆਕਾਰਾਂ, ਮਾਪਾਂ ਅਤੇ ਖਲਾਅ ਦਾ ਅਧਿਐਨ ਕਰਦੀ ਹੈ।}
\dictentrytr{gilt}{A bond issued by the UK government to borrow money.}{ਗਿਲਟ}{ਯੂਕੇ ਸਰਕਾਰ ਵੱਲੋਂ ਪੈਸੇ ਉਧਾਰ ਲੈਣ ਲਈ ਜਾਰੀ ਕੀਤਾ ਗਿਆ ਬਾਂਡ।}
\dictentrytr{gradient}{how steep a line is, found by how far it rises for each step across}{ਢਾਲ}{ਰੇਖਾ ਕਿੰਨੀ ਢਲਵੀਂ ਹੈ, ਇਹ ਦੇਖ ਕੇ ਪਤਾ ਲੱਗਦਾ ਹੈ ਕਿ ਉਹ ਹਰ ਇੱਕ ਕਦਮ ਪਾਰ ਕਰਨ 'ਤੇ ਕਿੰਨੀ ਉੱਪਰ ਚੜ੍ਹਦੀ ਹੈ।}
\dictentrytr{height}{the distance straight up from the base of a shape to its top, used to find area.}{ਉਚਾਈ}{ਕਿਸੇ ਆਕਾਰ ਦੇ ਅਧਾਰ ਤੋਂ ਉਸਦੇ ਸਿਖਰ ਤੱਕ ਸਿੱਧੀ ਉੱਪਰ ਵੱਲ ਦੂਰੀ, ਜੋ ਖੇਤਰਫਲ ਲੱਭਣ ਲਈ ਵਰਤੀ ਜਾਂਦੀ ਹੈ।}
\dictentrytr{hexagon}{a flat shape with six straight sides}{ਸ਼ਸ਼ਭੁਜ}{ਛੇ ਸਿੱਧੀਆਂ ਭੁਜਾਵਾਂ ਵਾਲੀ ਸਮਤਲ ਆਕਾਰ।}
\dictentrytr{highest common factor}{the largest whole number that divides exactly into two or more numbers}{ਮਹਾਨਤਮ ਸਾਂਝਾ ਗੁਣਨਖੰਡ}{ਸਭ ਤੋਂ ਵੱਡੀ ਪੂਰਨ ਸੰਖਿਆ ਜੋ ਦੋ ਜਾਂ ਵੱਧ ਸੰਖਿਆਵਾਂ ਨੂੰ ਪੂਰਾ-ਪੂਰਾ ਵੰਡਦੀ ਹੈ।}
\dictentrytr{hundredths}{the value of one part when a whole is divided into one hundred equal parts.}{ਸੌਵੇਂ}{ਇੱਕ ਪੂਰੇ ਨੂੰ ਸੌ ਬਰਾਬਰ ਹਿੱਸਿਆਂ ਵਿੱਚ ਵੰਡਣ 'ਤੇ ਇੱਕ ਹਿੱਸੇ ਦਾ ਮੁੱਲ।}
\dictentrytr{hypotenuse}{the longest side of a right-angled triangle, opposite the right angle}{ਕਰਨ}{ਸਮਕੋਣ ਤ੍ਰਿਭੁਜ ਦੀ ਸਭ ਤੋਂ ਲੰਬੀ ਭੁਜਾ, ਜੋ ਸਮਕੋਣ ਦੇ ਸਾਹਮਣੇ ਹੁੰਦੀ ਹੈ}
\dictentrytr{identity}{a statement that is true for every value of the letters in it}{ਸਰਵਸਮਿਕਾ}{ਇੱਕ ਬਿਆਨ ਜੋ ਉਸ ਵਿੱਚ ਆਏ ਅੱਖਰਾਂ ਦੇ ਹਰੇਕ ਮੁੱਲ ਲਈ ਸੱਚ ਹੁੰਦਾ ਹੈ।}
\dictentrytr{improper fraction}{a fraction with a top number larger than or equal to its bottom number}{ਅਨੁਚਿਤ ਭਿੰਨ}{ਉਹ ਭਿੰਨ ਜਿਸਦਾ ਉੱਪਰ ਵਾਲਾ ਅੰਕ ਆਪਣੇ ਹੇਠਲੇ ਅੰਕ ਤੋਂ ਵੱਡਾ ਜਾਂ ਬਰਾਬਰ ਹੋਵੇ।}
\dictentrytr{index}{the small raised number showing what power something is raised to}{ਘਾਤਾਂਕ}{ਉੱਪਰ ਲੱਗੀ ਛੋਟੀ ਸੰਖਿਆ ਜੋ ਦੱਸਦੀ ਹੈ ਕਿ ਕਿਸੇ ਚੀਜ਼ ਨੂੰ ਕਿਹੜੀ ਘਾਤ ਦਿੱਤੀ ਗਈ ਹੈ।}
\dictentrytr{inequality}{a statement that one quantity is larger or smaller than another}{ਅਸਮਾਨਤਾ}{ਇੱਕ ਬਿਆਨ ਜੋ ਦੱਸਦਾ ਹੈ ਕਿ ਇੱਕ ਮਾਤਰਾ ਦੂਜੀ ਨਾਲੋਂ ਵੱਡੀ ਜਾਂ ਛੋਟੀ ਹੈ।}
\dictentrytr{integer}{a whole number, which may be positive, negative or zero}{ਪੂਰਨ ਅੰਕ}{ਇੱਕ ਪੂਰਨ ਅੰਕ, ਜੋ ਸਕਾਰਾਤਮਕ, ਨਕਾਰਾਤਮਕ ਜਾਂ ਜ਼ੀਰੋ ਹੋ ਸਕਦਾ ਹੈ।}
\dictentrytr{intercept}{the point where a graph crosses an axis}{ਇੰਟਰਸੈਪਟ}{ਉਹ ਬਿੰਦੂ ਜਿੱਥੇ ਗ੍ਰਾਫ਼ ਇੱਕ ਧੁਰੇ ਨੂੰ ਕੱਟਦਾ ਹੈ}
\dictentrytr{interest}{Money paid for borrowing or earned on savings.}{ਵਿਆਜ}{ਉਧਾਰ ਲੈਣ ਲਈ ਦਿੱਤੇ ਜਾਂ ਬੱਚਤਾਂ 'ਤੇ ਕਮਾਏ ਗਏ ਪੈਸੇ।}
\dictentrytr{interest rate}{the percentage of an amount that is paid or earned as interest.}{ਵਿਆਜ ਦਰ}{ਕਿਸੇ ਰਕਮ ਦਾ ਉਹ ਪ੍ਰਤੀਸ਼ਤ ਜੋ ਵਿਆਜ ਦੇ ਰੂਪ ਵਿੱਚ ਦਿੱਤਾ ਜਾਂ ਕਮਾਇਆ ਜਾਂਦਾ ਹੈ।}
\dictentrytr{intersect}{to cross each other}{ਕੱਟਣਾ}{ਇੱਕ ਦੂਜੇ ਨੂੰ ਪਾਰ ਕਰਨਾ।}
\dictentrytr{inverse sine}{The operation that finds which corner has a given sine value.}{ਇਨਵਰਸ ਸਾਈਨ}{ਉਹ ਕਿਰਿਆ ਜੋ ਲੱਭਦੀ ਹੈ ਕਿ ਕਿਸ ਕੋਣ ਦਾ ਸਾਈਨ ਮੁੱਲ ਦਿੱਤਾ ਗਿਆ ਹੈ।}
\dictentrytr{isosceles triangle}{a three-sided shape with two sides equal in length}{ਸਮਦੋਭੁਜੀ ਤ੍ਰਿਭੁਜ}{ਇੱਕ ਤਿੰਨ-ਭੁਜਾਵਾਂ ਵਾਲੀ ਆਕਾਰ ਜਿਸਦੀਆਂ ਦੋ ਭੁਜਾਵਾਂ ਬਰਾਬਰ ਲੰਬਾਈ ਦੀਆਂ ਹੋਣ।}
\dictentrytr{iteration}{one complete repeat of a repeating process}{ਦੁਹਰਾਅ}{ਦੁਹਰਾਉਣ ਵਾਲੀ ਪ੍ਰਕਿਰਿਆ ਦਾ ਇੱਕ ਪੂਰਾ ਦੁਹਰਾਅ।}
\dictentrytr{iterative}{using a repeated rule, where each result is used to find the next one.}{ਦੁਹਰਾਊ}{ਦੁਹਰਾਏ ਗਏ ਨਿਯਮ ਦੀ ਵਰਤੋਂ ਕਰਦੇ ਹੋਏ, ਜਿੱਥੇ ਹਰ ਨਤੀਜਾ ਅਗਲਾ ਲੱਭਣ ਲਈ ਵਰਤਿਆ ਜਾਂਦਾ ਹੈ।}
\dictentrytr{iterative process}{a method that repeats a rule, using the previous answer to produce the next}{ਦੁਹਰਾਊ ਪ੍ਰਕਿਰਿਆ}{ਇੱਕ ਵਿਧੀ ਜੋ ਕਿਸੇ ਨਿਯਮ ਨੂੰ ਦੁਹਰਾਉਂਦੀ ਹੈ, ਪਿਛਲੇ ਜਵਾਬ ਨੂੰ ਅਗਲਾ ਜਵਾਬ ਪੈਦਾ ਕਰਨ ਲਈ ਵਰਤਦੀ ਹੈ।}
\dictentrytr{ladder}{an investment strategy that buys bonds with different maturity dates so payments arrive at regular times.}{ਲੈਡਰ}{ਇੱਕ ਨਿਵੇਸ਼ ਰਣਨੀਤੀ ਜੋ ਵੱਖ-ਵੱਖ ਪਰਿਪੱਕਤਾ ਮਿਤੀਆਂ ਵਾਲੇ ਬਾਂਡ ਖਰੀਦਦੀ ਹੈ ਤਾਂ ਜੋ ਭੁਗਤਾਨ ਨਿਯਮਤ ਸਮਿਆਂ 'ਤੇ ਆਉਣ।}
\dictentrytr{lattice method}{a written multiplication method that uses a grid to organise the calculation.}{ਜਾਲੀ ਵਿਧੀ}{ਗੁਣਾ ਕਰਨ ਦੀ ਇੱਕ ਲਿਖਤੀ ਵਿਧੀ ਜੋ ਗਣਨਾ ਨੂੰ ਵਿਵਸਥਿਤ ਕਰਨ ਲਈ ਗਰਿੱਡ ਦੀ ਵਰਤੋਂ ਕਰਦੀ ਹੈ।}
\dictentrytr{leg}{one of the two shorter sides of a right-angled triangle.}{ਲੇਗ}{ਸਮਕੋਣ ਤਿਕੋਣ ਦੀਆਂ ਦੋ ਛੋਟੀਆਂ ਭੁਜਾਵਾਂ ਵਿੱਚੋਂ ਇੱਕ।}
\dictentrytr{like terms}{terms that have the same letter or letters, or no letter, so they can be added or subtracted together.}{ਸਮਾਨ ਪਦ}{ਉਹ ਪਦ ਜਿਨ੍ਹਾਂ ਵਿੱਚ ਇੱਕੋ ਅੱਖਰ ਜਾਂ ਅੱਖਰ ਹੋਣ, ਜਾਂ ਕੋਈ ਅੱਖਰ ਨਾ ਹੋਵੇ, ਤਾਂ ਜੋ ਉਹਨਾਂ ਨੂੰ ਇਕੱਠੇ ਜੋੜਿਆ ਜਾਂ ਘਟਾਇਆ ਜਾ ਸਕੇ।}
\dictentrytr{line segment}{a straight part of a line that starts at one point and ends at another.}{ਰੇਖਾ ਖੰਡ}{ਰੇਖਾ ਦਾ ਸਿੱਧਾ ਹਿੱਸਾ ਜੋ ਇੱਕ ਬਿੰਦੂ ਤੋਂ ਸ਼ੁਰੂ ਹੋ ਕੇ ਦੂਜੇ ਬਿੰਦੂ 'ਤੇ ਖਤਮ ਹੁੰਦਾ ਹੈ।}
\dictentrytr{linear}{making a straight line when drawn as a graph}{ਰੇਖੀ}{ਜਦੋਂ ਗ੍ਰਾਫ਼ ਵਜੋਂ ਖਿੱਚਿਆ ਜਾਵੇ ਤਾਂ ਸਿੱਧੀ ਰੇਖਾ ਬਣਾਉਣ ਵਾਲਾ}
\dictentrytr{lowest common multiple}{the smallest number that two or more numbers all divide into exactly}{ਲਘੁਤਮ ਸਾਂਝਾ ਗੁਣਜ}{ਸਭ ਤੋਂ ਛੋਟੀ ਸੰਖਿਆ ਜਿਸ ਨੂੰ ਦੋ ਜਾਂ ਵੱਧ ਸੰਖਿਆਵਾਂ ਪੂਰਾ-ਪੂਰਾ ਵੰਡਦੀਆਂ ਹਨ।}
\dictentrytr{magnitude}{the size of a number without its sign.}{ਪਰਿਮਾਣ}{ਚਿੰਨ੍ਹ ਤੋਂ ਬਿਨਾਂ ਕਿਸੇ ਸੰਖਿਆ ਦਾ ਆਕਾਰ।}
\dictentrytr{major arc}{the longer of the two parts of the circumference between two points on a circle.}{ਵੱਡਾ ਚਾਪ}{ਕਿਸੇ ਚੱਕਰ ਉੱਤੇ ਦੋ ਬਿੰਦੂਆਂ ਦੇ ਵਿਚਕਾਰ ਪਰਿਧੀ ਦੇ ਦੋ ਹਿੱਸਿਆਂ ਵਿੱਚੋਂ ਲੰਬਾ ਹਿੱਸਾ।}
\dictentrytr{major segment}{the larger of the two parts of a circle cut off by a chord.}{ਵੱਡਾ ਚੱਕਰ ਖੰਡ}{ਕਿਸੇ ਜੀਵਾ ਦੁਆਰਾ ਚੱਕਰ ਦੇ ਕੱਟੇ ਗਏ ਦੋ ਹਿੱਸਿਆਂ ਵਿੱਚੋਂ ਵੱਡਾ ਹਿੱਸਾ।}
\dictentrytr{mapping}{matching or sending each point of one shape to a point of another.}{ਮੈਪਿੰਗ}{ਇੱਕ ਆਕਾਰ ਦੇ ਹਰੇਕ ਬਿੰਦੂ ਨੂੰ ਦੂਜੇ ਆਕਾਰ ਦੇ ਕਿਸੇ ਬਿੰਦੂ ਨਾਲ ਮਿਲਾਉਣਾ ਜਾਂ ਭੇਜਣਾ।}
\dictentrytr{mass}{the amount of matter in an object, measured in grams or kilograms.}{ਪੁੰਜ}{ਕਿਸੇ ਵਸਤੂ ਵਿੱਚ ਪਦਾਰਥ ਦੀ ਮਾਤਰਾ, ਜੋ ਗ੍ਰਾਮ ਜਾਂ ਕਿਲੋਗ੍ਰਾਮ ਵਿੱਚ ਮਾਪੀ ਜਾਂਦੀ ਹੈ।}
\dictentrytr{maturity}{the time when a bond is paid back.}{ਪਰਿਪੱਕਤੀ}{ਉਹ ਸਮਾਂ ਜਦੋਂ ਬਾਂਡ ਦਾ ਪੈਸਾ ਵਾਪਸ ਕੀਤਾ ਜਾਂਦਾ ਹੈ।}
\dictentrytr{mean}{the average found by adding all the values and dividing by how many there are}{ਔਸਤ}{ਔਸਤ ਜੋ ਸਾਰੇ ਮੁੱਲਾਂ ਨੂੰ ਜੋੜ ਕੇ ਅਤੇ ਉਹਨਾਂ ਦੀ ਗਿਣਤੀ ਨਾਲ ਭਾਗ ਕਰਕੇ ਪ੍ਰਾਪਤ ਕੀਤੀ ਜਾਂਦੀ ਹੈ}
\dictentrytr{median}{the middle value when the data is put in order}{ਮੱਧਿਕਾ}{ਵਿਚਕਾਰਲਾ ਮੁੱਲ ਜਦੋਂ ਅੰਕੜਿਆਂ ਨੂੰ ਕ੍ਰਮ ਵਿੱਚ ਰੱਖਿਆ ਜਾਂਦਾ ਹੈ}
\dictentrytr{minor arc}{the shorter of the two parts of the circumference between two points on a circle.}{ਛੋਟਾ ਚਾਪ}{ਕਿਸੇ ਚੱਕਰ ਉੱਤੇ ਦੋ ਬਿੰਦੂਆਂ ਦੇ ਵਿਚਕਾਰ ਪਰਿਧੀ ਦੇ ਦੋ ਹਿੱਸਿਆਂ ਵਿੱਚੋਂ ਛੋਟਾ ਹਿੱਸਾ।}
\dictentrytr{minor segment}{the smaller of the two parts of a circle cut off by a chord.}{ਛੋਟਾ ਚੱਕਰ ਖੰਡ}{ਕਿਸੇ ਜੀਵਾ ਦੁਆਰਾ ਚੱਕਰ ਦੇ ਕੱਟੇ ਗਏ ਦੋ ਹਿੱਸਿਆਂ ਵਿੱਚੋਂ ਛੋਟਾ ਹਿੱਸਾ।}
\dictentrytr{mixed number}{a number made of a whole amount and a fraction, such as 3 1/2}{ਮਿਸ਼ਰਤ ਸੰਖਿਆ}{ਇੱਕ ਸੰਖਿਆ ਜੋ ਪੂਰੀ ਮਾਤਰਾ ਅਤੇ ਭਿੰਨ ਤੋਂ ਬਣੀ ਹੋਵੇ, ਜਿਵੇਂ 3 1/2।}
\dictentrytr{mode}{the value that appears most often}{ਬਹੁਲਕ}{ਉਹ ਮੁੱਲ ਜੋ ਸਭ ਤੋਂ ਵੱਧ ਵਾਰ ਆਉਂਦਾ ਹੈ}
\dictentrytr{multiple}{a number you get by multiplying another number by a whole number}{ਗੁਣਜ}{ਇੱਕ ਸੰਖਿਆ ਜੋ ਕਿਸੇ ਹੋਰ ਸੰਖਿਆ ਨੂੰ ਇੱਕ ਪੂਰਨ ਅੰਕ ਨਾਲ ਗੁਣਾ ਕਰਨ 'ਤੇ ਮਿਲਦੀ ਹੈ।}
\dictentrytr{multiply}{to add a number to itself a given number of times}{ਗੁਣਾ ਕਰਨਾ}{ਕਿਸੇ ਸੰਖਿਆ ਨੂੰ ਆਪਣੇ ਆਪ ਵਿੱਚ ਇੱਕ ਨਿਸ਼ਚਿਤ ਗਿਣਤੀ ਵਾਰ ਜੋੜਨਾ।}
\dictentrytr{negative}{below zero; a number or term with a minus sign is below zero.}{਋ਣ}{ਸਿਫ਼ਰ ਤੋਂ ਹੇਠਾਂ; ਘਟਾਓ ਦੇ ਚਿੰਨ੍ਹ ਵਾਲੀ ਸੰਖਿਆ ਜਾਂ ਪਦ ਸਿਫ਼ਰ ਤੋਂ ਹੇਠਾਂ ਹੁੰਦਾ ਹੈ।}
\dictentrytr{negative number}{a number less than zero, written with a minus sign.}{਋ਣ ਸੰਖਿਆ}{ਇੱਕ ਸੰਖਿਆ ਜੋ ਸਿਫ਼ਰ ਤੋਂ ਘੱਟ ਹੋਵੇ, ਘਟਾਓ ਦੇ ਚਿੰਨ੍ਹ ਨਾਲ ਲਿਖੀ ਜਾਂਦੀ ਹੈ।}
\dictentrytr{non-reflex angle}{the smaller angle between two lines, measuring less than 180 degrees}{ਗੈਰ-ਰਿਫਲੈਕਸ ਕੋਣ}{ਦੋ ਰੇਖਾਵਾਂ ਦੇ ਵਿਚਕਾਰ ਬਣਿਆ ਛੋਟਾ ਕੋਣ, ਜੋ 180 ਡਿਗਰੀ ਤੋਂ ਘੱਟ ਹੁੰਦਾ ਹੈ।}
\dictentrytr{nth term}{a rule that gives any number in a sequence from its position}{nਵਾਂ ਪਦ}{ਇੱਕ ਨਿਯਮ ਜੋ ਲੜੀ ਵਿੱਚ ਕਿਸੇ ਵੀ ਸੰਖਿਆ ਨੂੰ ਉਸਦੀ ਸਥਿਤੀ ਤੋਂ ਦਿੰਦਾ ਹੈ}
\dictentrytr{number line}{a straight line with numbers marked at equal distances, used to show value and order.}{ਸੰਖਿਆ ਰੇਖਾ}{ਇੱਕ ਸਿੱਧੀ ਰੇਖਾ ਜਿਸ ਉੱਤੇ ਸੰਖਿਆਵਾਂ ਬਰਾਬਰ ਦੂਰੀ 'ਤੇ ਅੰਕਿਤ ਹੁੰਦੀਆਂ ਹਨ, ਮੁੱਲ ਅਤੇ ਕ੍ਰਮ ਦਿਖਾਉਣ ਲਈ ਵਰਤੀ ਜਾਂਦੀ ਹੈ।}
\dictentrytr{numerator}{the number above the line in a fraction, showing how many equal parts are taken}{ਅੰਸ਼}{ਭਿੰਨ ਵਿੱਚ ਰੇਖਾ ਦੇ ਉੱਪਰ ਵਾਲੀ ਸੰਖਿਆ, ਜੋ ਦੱਸਦੀ ਹੈ ਕਿ ਕਿੰਨੇ ਬਰਾਬਰ ਹਿੱਸੇ ਲਏ ਗਏ ਹਨ।}
\dictentrytr{odd}{a whole number that cannot be split exactly into two equal whole numbers.}{ਵਿਸਮ}{ਇੱਕ ਪੂਰਨ ਸੰਖਿਆ ਜਿਸਨੂੰ ਦੋ ਬਰਾਬਰ ਪੂਰਨ ਸੰਖਿਆਵਾਂ ਵਿੱਚ ਬਿਲਕੁਲ ਵੰਡਿਆ ਨਹੀਂ ਜਾ ਸਕਦਾ।}
\dictentrytr{odd number}{a whole number that cannot be divided exactly by two.}{ਟਾਂਕ ਸੰਖਿਆ}{ਉਹ ਪੂਰਨ ਸੰਖਿਆ ਜਿਸਨੂੰ ਦੋ ਨਾਲ ਪੂਰੀ ਤਰ੍ਹਾਂ ਵੰਡਿਆ ਨਹੀਂ ਜਾ ਸਕਦਾ।}
\dictentrytr{of}{in fraction work, it means multiply: taking a fraction of an amount uses multiplication}{ਦਾ}{ਭਿੰਨ ਦੇ ਕੰਮ ਵਿੱਚ, ਇਸਦਾ ਅਰਥ ਹੈ ਗੁਣਾ ਕਰਨਾ: ਕਿਸੇ ਰਕਮ ਦਾ ਭਿੰਨ ਕੱਢਣ ਵਿੱਚ ਗੁਣਾ ਵਰਤਿਆ ਜਾਂਦਾ ਹੈ।}
\dictentrytr{operation}{Something done to a number, such as adding, subtracting, multiplying or dividing.}{ਕਿਰਿਆ}{ਕਿਸੇ ਸੰਖਿਆ ਨਾਲ ਕੀਤੀ ਜਾਣ ਵਾਲੀ ਚੀਜ਼, ਜਿਵੇਂ ਜੋੜਨਾ, ਘਟਾਉਣਾ, ਗੁਣਾ ਕਰਨਾ ਜਾਂ ਭਾਗ ਕਰਨਾ।}
\dictentrytr{opposite}{in a three-sided shape with a right angle, the side across from a chosen corner.}{ਸਾਹਮਣੇ ਵਾਲਾ}{ਸਮਕੋਣ ਵਾਲੇ ਤਿੰਨ-ਭੁਜਾ ਆਕਾਰ ਵਿੱਚ, ਕਿਸੇ ਚੁਣੇ ਹੋਏ ਕੋਨੇ ਦੇ ਸਾਹਮਣੇ ਵਾਲੀ ਭੁਜਾ।}
\dictentrytr{order}{to arrange numbers from smallest to largest, or from largest to smallest.}{ਕ੍ਰਮਬੱਧ ਕਰਨਾ}{ਸੰਖਿਆਵਾਂ ਨੂੰ ਸਭ ਤੋਂ ਛੋਟੇ ਤੋਂ ਸਭ ਤੋਂ ਵੱਡੇ, ਜਾਂ ਸਭ ਤੋਂ ਵੱਡੇ ਤੋਂ ਸਭ ਤੋਂ ਛੋਟੇ ਕ੍ਰਮ ਵਿੱਚ ਲਗਾਉਣਾ।}
\dictentrytr{origin}{the fixed point (0,0) on a coordinate grid.}{ਮੂਲ ਬਿੰਦੂ}{ਨਿਰਦੇਸ਼ਾਂਕ ਗਰਿੱਡ ਉੱਤੇ ਸਥਿਰ ਬਿੰਦੂ (0,0)।}
\dictentrytr{outlier}{a value far away from all the others}{ਬਾਹਰੀ ਮੁੱਲ}{ਇੱਕ ਮੁੱਲ ਜੋ ਬਾਕੀ ਸਾਰੇ ਮੁੱਲਾਂ ਤੋਂ ਬਹੁਤ ਦੂਰ ਹੁੰਦਾ ਹੈ}
\dictentrytr{parallel}{always the same distance apart and never meeting}{ਸਮਾਂਤਰ}{ਹਮੇਸ਼ਾ ਇੱਕੋ ਦੂਰੀ ਉੱਤੇ ਰਹਿਣ ਵਾਲੇ ਅਤੇ ਕਦੇ ਨਾ ਮਿਲਣ ਵਾਲੇ}
\dictentrytr{parentheses}{the round brackets ( ) used in calculations to show which part to do first.}{ਕੋਸ਼ਟਕ}{ਗੋਲ ਬਰੈਕਟਾਂ ( ) ਜੋ ਹਿਸਾਬਾਂ ਵਿੱਚ ਇਹ ਦਿਖਾਉਣ ਲਈ ਵਰਤੀਆਂ ਜਾਂਦੀਆਂ ਹਨ ਕਿ ਕਿਹੜਾ ਹਿੱਸਾ ਪਹਿਲਾਂ ਕਰਨਾ ਹੈ।}
\dictentrytr{per cent}{for every hundred; a number out of one hundred.}{ਪ੍ਰਤੀਸ਼ਤ}{ਹਰ ਸੌ ਲਈ; ਸੌ ਵਿੱਚੋਂ ਇੱਕ ਸੰਖਿਆ।}
\dictentrytr{percent}{for every hundred; shown with a \% sign.}{ਪ੍ਰਤੀਸ਼ਤ}{ਹਰ ਸੌ ਲਈ; \% ਚਿੰਨ੍ਹ ਨਾਲ ਦਿਖਾਇਆ ਜਾਂਦਾ ਹੈ।}
\dictentrytr{percentage}{a number of parts out of a hundred}{ਪ੍ਰਤੀਸ਼ਤ}{ਸੌ ਵਿੱਚੋਂ ਭਾਗਾਂ ਦੀ ਗਿਣਤੀ।}
\dictentrytr{perimeter}{the total distance around the edge of a shape}{ਪਰਿਮਾਪ}{ਕਿਸੇ ਆਕਾਰ ਦੇ ਕਿਨਾਰੇ ਦੇ ਚਾਰੇ ਪਾਸੇ ਦੀ ਕੁੱਲ ਦੂਰੀ}
\dictentrytr{perpendicular}{meeting at a right angle}{ਲੰਬਵਤ}{ਸਮਕੋਣ ਉੱਤੇ ਮਿਲਣ ਵਾਲੇ}
\dictentrytr{pi}{the special number about 3.14 used when working with circles}{ਪਾਈ}{ਵਿਸ਼ੇਸ਼ ਸੰਖਿਆ ਲਗਭਗ 3.14 ਜੋ ਚੱਕਰਾਂ ਨਾਲ ਕੰਮ ਕਰਨ ਵੇਲੇ ਵਰਤੀ ਜਾਂਦੀ ਹੈ।}
\dictentrytr{pie chart}{a circular chart divided into slices to show parts of a whole.}{ਪਾਈ ਚਾਰਟ}{ਇੱਕ ਗੋਲ ਚਾਰਟ ਜਿਸਨੂੰ ਟੁਕੜਿਆਂ ਵਿੱਚ ਵੰਡ ਕੇ ਪੂਰੇ ਦੇ ਹਿੱਸੇ ਦਿਖਾਏ ਜਾਂਦੇ ਹਨ।}
\dictentrytr{place value}{the value a digit has because of where it sits in a number}{ਸਥਾਨ ਮੁੱਲ}{ਇੱਕ ਅੰਕ ਦਾ ਮੁੱਲ ਜੋ ਸੰਖਿਆ ਵਿੱਚ ਉਸਦੀ ਸਥਿਤੀ ਕਾਰਨ ਹੁੰਦਾ ਹੈ।}
\dictentrytr{polygon}{a flat shape with straight sides}{ਬਹੁਭੁਜ}{ਸਿੱਧੀਆਂ ਭੁਜਾਵਾਂ ਵਾਲੀ ਸਮਤਲ ਸ਼ਕਲ}
\dictentrytr{position-to-term rule}{a way to calculate any term in a sequence from its position, such as 3n+1.}{ਸਥਿਤੀ-ਤੋਂ-ਪਦ ਨਿਯਮ}{ਕਿਸੇ ਲੜੀ ਵਿੱਚ ਕਿਸੇ ਵੀ ਪਦ ਨੂੰ ਉਸਦੀ ਸਥਿਤੀ ਤੋਂ ਗਿਣਨ ਦਾ ਤਰੀਕਾ, ਜਿਵੇਂ 3n+1।}
\dictentrytr{positive}{a number greater than zero.}{ਧਨਾਤਮਕ}{ਸਿਫ਼ਰ ਤੋਂ ਵੱਡੀ ਸੰਖਿਆ।}
\dictentrytr{power}{how many times a number is multiplied by itself}{ਘਾਤ}{ਇੱਕ ਸੰਖਿਆ ਨੂੰ ਆਪਣੇ ਆਪ ਨਾਲ ਕਿੰਨੀ ਵਾਰ ਗੁਣਾ ਕੀਤਾ ਜਾਂਦਾ ਹੈ।}
\dictentrytr{prime}{a whole number greater than 1 that divides only by itself and 1}{ਅਭਾਜ ਸੰਖਿਆ}{1 ਤੋਂ ਵੱਡੀ ਪੂਰਨ ਸੰਖਿਆ ਜੋ ਸਿਰਫ ਆਪਣੇ ਆਪ ਅਤੇ 1 ਨਾਲ ਵੰਡੀ ਜਾਂਦੀ ਹੈ।}
\dictentrytr{prime number}{a whole number greater than 1 whose only factor is 1 and itself.}{ਅਭਾਜ ਸੰਖਿਆ}{1 ਤੋਂ ਵੱਡੀ ਪੂਰਨ ਸੰਖਿਆ ਜਿਸਦਾ ਇੱਕੋ ਇੱਕ ਗੁਣਨਖੰਡ 1 ਅਤੇ ਉਹ ਖੁਦ ਹੈ।}
\dictentrytr{principal}{The original amount of money borrowed or lent, before interest or extra payments.}{ਮੂਲ ਰਕਮ}{ਉਧਾਰ ਲਈ ਜਾਂ ਦਿੱਤੀ ਗਈ ਮੂਲ ਰਕਮ, ਵਿਆਜ ਜਾਂ ਵਾਧੂ ਭੁਗਤਾਨਾਂ ਤੋਂ ਪਹਿਲਾਂ।}
\dictentrytr{probability}{how likely something is, given as a number between 0 and 1}{ਸੰਭਾਵਨਾ}{ਕਿਸੇ ਘਟਨਾ ਦੀ ਸੰਭਾਵਨਾ, ਜੋ 0 ਅਤੇ 1 ਦੇ ਵਿਚਕਾਰ ਇੱਕ ਸੰਖਿਆ ਦੇ ਰੂਪ ਵਿੱਚ ਦਿੱਤੀ ਜਾਂਦੀ ਹੈ}
\dictentrytr{product}{the answer when two or more numbers are multiplied}{ਗੁਣਨਫਲ}{ਦੋ ਜਾਂ ਵੱਧ ਸੰਖਿਆਵਾਂ ਨੂੰ ਗੁਣਾ ਕਰਨ 'ਤੇ ਮਿਲਣ ਵਾਲਾ ਉੱਤਰ।}
\dictentrytr{proof}{a logical argument showing that a mathematical statement is always true.}{ਉਪਪੱਤੀ}{ਇੱਕ ਤਰਕਪੂਰਨ ਦਲੀਲ ਜੋ ਦਿਖਾਉਂਦੀ ਹੈ ਕਿ ਇੱਕ ਗਣਿਤਕ ਕਥਨ ਹਮੇਸ਼ਾ ਸੱਚ ਹੁੰਦਾ ਹੈ।}
\dictentrytr{proper fraction}{a fraction with a top number smaller than its bottom number}{ਉਚਿਤ ਭਿੰਨ}{ਇੱਕ ਭਿੰਨ ਜਿਸਦਾ ਉੱਪਰਲਾ ਅੰਕ ਉਸਦੇ ਹੇਠਲੇ ਅੰਕ ਤੋਂ ਛੋਟਾ ਹੁੰਦਾ ਹੈ।}
\dictentrytr{proportion}{a relationship where two amounts change by the same factor}{ਸਮਾਨੁਪਾਤ}{ਇੱਕ ਸਬੰਧ ਜਿਸ ਵਿੱਚ ਦੋ ਮਾਤਰਾਵਾਂ ਇੱਕੋ ਗੁਣਨਖੰਡ ਨਾਲ ਬਦਲਦੀਆਂ ਹਨ।}
\dictentrytr{protractor}{a tool used to measure angles.}{ਕੋਣਮਾਪਕ}{ਕੋਣਾਂ ਨੂੰ ਨਾਪਣ ਲਈ ਵਰਤਿਆ ਜਾਣ ਵਾਲਾ ਸੰਦ।}
\dictentrytr{pythagoras}{a rule for finding a missing side of a right-angled triangle.}{ਪਾਇਥਾਗੋਰਸ}{ਸਮਕੋਣ ਤਿਕੋਣ ਦੀ ਲਾਪਤਾ ਭੁਜਾ ਲੱਭਣ ਲਈ ਇੱਕ ਨਿਯਮ।}
\dictentrytr{pythagoras' theorem}{in a right-angled triangle, the longest side times itself equals the other two sides times themselves, added together.}{ਪਾਇਥਾਗੋਰਸ ਦਾ ਸਿਧਾਂਤ}{ਸਮਕੋਣ ਤਿਕੋਣ ਵਿੱਚ, ਸਭ ਤੋਂ ਲੰਬੀ ਭੁਜਾ ਦਾ ਆਪਣੇ ਨਾਲ ਗੁਣਾ, ਬਾਕੀ ਦੋ ਭੁਜਾਵਾਂ ਦੇ ਆਪਣੇ ਨਾਲ ਗੁਣਾ ਦੇ ਜੋੜ ਦੇ ਬਰਾਬਰ ਹੁੰਦਾ ਹੈ।}
\dictentrytr{quadratic}{an expression or equation whose highest power is a square}{ਦੋਘਾਤੀ}{ਇੱਕ ਵਿਅੰਜਕ ਜਾਂ ਸਮੀਕਰਨ ਜਿਸਦੀ ਸਭ ਤੋਂ ਉੱਚੀ ਘਾਤ ਵਰਗ ਹੋਵੇ}
\dictentrytr{quarter}{one of four equal parts of one.}{ਚੌਥਾਈ}{ਇੱਕ ਦੇ ਚਾਰ ਬਰਾਬਰ ਹਿੱਸਿਆਂ ਵਿੱਚੋਂ ਇੱਕ।}
\dictentrytr{quartile}{a value marking one quarter of the way through ordered data}{ਚਤੁਰਥਾਂਸ਼}{ਇੱਕ ਮੁੱਲ ਜੋ ਕ੍ਰਮਬੱਧ ਅੰਕੜਿਆਂ ਦਾ ਇੱਕ ਚੌਥਾਈ ਹਿੱਸਾ ਦਰਸਾਉਂਦਾ ਹੈ}
\dictentrytr{quotient}{the answer when one number is divided by another}{ਭਾਗਫਲ}{ਇੱਕ ਸੰਖਿਆ ਨੂੰ ਦੂਜੀ ਨਾਲ ਵੰਡਣ 'ਤੇ ਮਿਲਣ ਵਾਲਾ ਉੱਤਰ।}
\dictentrytr{radius}{the distance from the centre of a circle to its edge}{ਅਰਧਵਿਆਸ}{ਚੱਕਰ ਦੇ ਕੇਂਦਰ ਤੋਂ ਉਸਦੇ ਕਿਨਾਰੇ ਤੱਕ ਦੀ ਦੂਰੀ}
\dictentrytr{range}{the difference between the largest and smallest values}{ਪਸਾਰ}{ਸਭ ਤੋਂ ਵੱਡੇ ਅਤੇ ਸਭ ਤੋਂ ਛੋਟੇ ਮੁੱਲ ਵਿਚਕਾਰ ਦਾ ਅੰਤਰ}
\dictentrytr{rate}{a comparison of two different kinds of amount, such as litres per minute.}{ਦਰ}{ਦੋ ਵੱਖ-ਵੱਖ ਕਿਸਮਾਂ ਦੀਆਂ ਮਾਤਰਾਵਾਂ ਦੀ ਤੁਲਨਾ, ਜਿਵੇਂ ਲਿਟਰ ਪ੍ਰਤੀ ਮਿੰਟ।}
\dictentrytr{rate of change}{how quickly a measurement increases or decreases over time.}{ਤਬਦੀਲੀ ਦੀ ਦਰ}{ਸਮੇਂ ਦੇ ਨਾਲ ਕੋਈ ਮਾਪ ਕਿੰਨੀ ਤੇਜ਼ੀ ਨਾਲ ਵਧਦਾ ਜਾਂ ਘਟਦਾ ਹੈ।}
\dictentrytr{ratio}{a way of comparing two or more amounts, showing how much of one there is for each of the other}{ਅਨੁਪਾਤ}{ਦੋ ਜਾਂ ਵੱਧ ਮਾਤਰਾਵਾਂ ਦੀ ਤੁਲਨਾ ਕਰਨ ਦਾ ਤਰੀਕਾ, ਜੋ ਦੱਸਦਾ ਹੈ ਕਿ ਦੂਜੀ ਦੇ ਮੁਕਾਬਲੇ ਇੱਕ ਕਿੰਨੀ ਹੈ।}
\dictentrytr{ray}{a straight line that starts at one point and goes on forever in one direction}{ਕਿਰਨ}{ਇੱਕ ਸਿੱਧੀ ਰੇਖਾ ਜੋ ਇੱਕ ਬਿੰਦੂ ਤੋਂ ਸ਼ੁਰੂ ਹੁੰਦੀ ਹੈ ਅਤੇ ਇੱਕ ਦਿਸ਼ਾ ਵਿੱਚ ਹਮੇਸ਼ਾ ਲਈ ਚੱਲਦੀ ਰਹਿੰਦੀ ਹੈ।}
\dictentrytr{rearrange}{to change the order of a formula so a different letter is on its own}{ਪੁਨਰ-ਵਿਵਸਥਿਤ ਕਰਨਾ}{ਕਿਸੇ ਸੂਤਰ ਦਾ ਕ੍ਰਮ ਬਦਲਣਾ ਤਾਂ ਜੋ ਕੋਈ ਹੋਰ ਅੱਖਰ ਇਕੱਲਾ ਹੋ ਜਾਵੇ}
\dictentrytr{reciprocal}{the number you multiply by to get 1, found by turning a fraction upside down}{ਉਲਟ}{ਉਹ ਸੰਖਿਆ ਜਿਸ ਨਾਲ ਗੁਣਾ ਕਰਨ 'ਤੇ 1 ਮਿਲਦਾ ਹੈ, ਭਿੰਨ ਨੂੰ ਉਲਟਾ ਕੇ ਲੱਭੀ ਜਾਂਦੀ ਹੈ।}
\dictentrytr{rectangle}{a flat shape with four straight sides and four square corners.}{ਆਇਤ}{ਇੱਕ ਸਮਤਲ ਆਕਾਰ ਜਿਸਦੀਆਂ ਚਾਰ ਸਿੱਧੀਆਂ ਭੁਜਾਵਾਂ ਅਤੇ ਚਾਰ ਸਮਕੋਣ ਹੋਣ।}
\dictentrytr{recur}{to repeat the same digits forever after the decimal point}{ਦੁਹਰਾਉਣਾ}{ਦਸ਼ਮਲਵ ਬਿੰਦੂ ਤੋਂ ਬਾਅਦ ਇੱਕੋ ਜਿਹੇ ਅੰਕਾਂ ਨੂੰ ਹਮੇਸ਼ਾ ਲਈ ਦੁਹਰਾਉਣਾ।}
\dictentrytr{recurrence}{a rule that defines each term using one or more earlier terms}{ਪੁਨਰਾਵਰਤਨ}{ਇੱਕ ਨਿਯਮ ਜੋ ਹਰ ਪਦ ਨੂੰ ਇੱਕ ਜਾਂ ਵੱਧ ਪਿਛਲੇ ਪਦਾਂ ਦੀ ਵਰਤੋਂ ਨਾਲ ਪਰਿਭਾਸ਼ਿਤ ਕਰਦਾ ਹੈ।}
\dictentrytr{recurring decimal}{a decimal in which a digit or group of digits repeats forever.}{ਆਵਰਤੀ ਦਸ਼ਮਲਵ}{ਇੱਕ ਦਸ਼ਮਲਵ ਜਿਸ ਵਿੱਚ ਕੋਈ ਅੰਕ ਜਾਂ ਅੰਕਾਂ ਦਾ ਸਮੂਹ ਹਮੇਸ਼ਾ ਦੁਹਰਾਇਆ ਜਾਂਦਾ ਹੈ।}
\dictentrytr{regular}{having all sides equal and all angles equal}{ਨਿਯਮਿਤ}{ਜਿਸਦੀਆਂ ਸਾਰੀਆਂ ਭੁਜਾਵਾਂ ਬਰਾਬਰ ਅਤੇ ਸਾਰੇ ਕੋਣ ਬਰਾਬਰ ਹੋਣ।}
\dictentrytr{remainder}{the amount left over when a number cannot be split into equal whole groups}{ਬਾਕੀ}{ਜਦੋਂ ਕਿਸੇ ਸੰਖਿਆ ਨੂੰ ਬਰਾਬਰ ਪੂਰਨ ਸਮੂਹਾਂ ਵਿੱਚ ਵੰਡਿਆ ਨਹੀਂ ਜਾ ਸਕਦਾ ਤਾਂ ਬਚੀ ਹੋਈ ਮਾਤਰਾ।}
\dictentrytr{resultant}{the single vector that has the same overall effect as adding two or more vectors.}{ਪਰਿਣਾਮੀ}{ਉਹ ਇੱਕੋ ਵੈਕਟਰ ਜਿਸਦਾ ਸਮੁੱਚਾ ਪ੍ਰਭਾਵ ਦੋ ਜਾਂ ਵੱਧ ਵੈਕਟਰਾਂ ਨੂੰ ਜੋੜਨ ਦੇ ਬਰਾਬਰ ਹੁੰਦਾ ਹੈ।}
\dictentrytr{return}{The amount of money gained or lost on an amount lent or invested.}{ਰਿਟਰਨ}{ਕਿਸੇ ਉਧਾਰ ਦਿੱਤੀ ਜਾਂ ਨਿਵੇਸ਼ ਕੀਤੀ ਗਈ ਰਕਮ 'ਤੇ ਕਮਾਈ ਜਾਂ ਗੁਆਈ ਗਈ ਰਕਮ।}
\dictentrytr{reverse}{The same vector with every number changed to the opposite sign, so it points back.}{ਉਲਟ}{ਉਹੀ ਵੈਕਟਰ ਜਿਸਦੀ ਹਰ ਸੰਖਿਆ ਦਾ ਚਿੰਨ੍ਹ ਉਲਟਾ ਦਿੱਤਾ ਗਿਆ ਹੋਵੇ, ਤਾਂ ਜੋ ਇਹ ਵਾਪਸ ਮੁੜੇ।}
\dictentrytr{right angle}{a square corner, exactly a quarter of a full turn.}{ਸਮਕੋਣ}{ਸਮਕੋਣ ਵਾਲਾ ਕੋਨਾ, ਜੋ ਪੂਰੇ ਮੋੜ ਦਾ ਬਿਲਕੁਲ ਚੌਥਾਈ ਹਿੱਸਾ ਹੁੰਦਾ ਹੈ।}
\dictentrytr{right triangle}{a triangle with one angle that is a quarter turn.}{ਸਮਕੋਣ ਤਿਕੋਣ}{ਇੱਕ ਤਿਕੋਣ ਜਿਸਦਾ ਇੱਕ ਕੋਣ ਚੌਥਾਈ ਮੋੜ ਦਾ ਹੁੰਦਾ ਹੈ।}
\dictentrytr{right-angled triangle}{a three-sided shape with one corner of 90 degrees.}{ਸਮਕੋਣ ਤਿਕੋਣ}{ਤਿੰਨ ਭੁਜਾਵਾਂ ਵਾਲੀ ਸ਼ਕਲ ਜਿਸਦਾ ਇੱਕ ਕੋਨਾ 90 ਡਿਗਰੀ ਦਾ ਹੁੰਦਾ ਹੈ।}
\dictentrytr{root}{a number that gives another number when multiplied by itself a given number of times}{ਮੂਲ}{ਇੱਕ ਸੰਖਿਆ ਜੋ ਆਪਣੇ ਆਪ ਨਾਲ ਇੱਕ ਨਿਸ਼ਚਿਤ ਗਿਣਤੀ ਵਾਰ ਗੁਣਾ ਹੋਣ 'ਤੇ ਕੋਈ ਹੋਰ ਸੰਖਿਆ ਦਿੰਦੀ ਹੈ।}
\dictentrytr{round}{to replace a number with a nearby simpler one}{ਗੋਲ ਕਰਨਾ}{ਕਿਸੇ ਸੰਖਿਆ ਨੂੰ ਨੇੜੇ ਦੀ ਸੌਖੀ ਸੰਖਿਆ ਨਾਲ ਬਦਲਣਾ।}
\dictentrytr{rule}{a mathematical instruction that works out every number in a sequence}{ਨਿਯਮ}{ਇੱਕ ਗਣਿਤ ਦਾ ਨਿਰਦੇਸ਼ ਜੋ ਲੜੀ ਵਿੱਚ ਹਰ ਸੰਖਿਆ ਕੱਢਦਾ ਹੈ।}
\dictentrytr{sample}{a smaller group chosen to stand for a larger one}{ਨਮੂਨਾ}{ਇੱਕ ਛੋਟਾ ਸਮੂਹ ਜੋ ਵੱਡੇ ਸਮੂਹ ਦੀ ਨੁਮਾਇੰਦਗੀ ਕਰਨ ਲਈ ਚੁਣਿਆ ਜਾਂਦਾ ਹੈ}
\dictentrytr{scale}{to enlarge or shrink a shape so every length changes in the same way}{ਸਕੇਲ}{ਕਿਸੇ ਆਕਾਰ ਨੂੰ ਵੱਡਾ ਜਾਂ ਛੋਟਾ ਕਰਨਾ ਤਾਂ ਜੋ ਹਰ ਲੰਬਾਈ ਇੱਕੋ ਤਰੀਕੇ ਨਾਲ ਬਦਲੇ।}
\dictentrytr{scale factor}{the number every length is multiplied by when a shape is enlarged}{ਸਕੇਲ ਗੁਣਕ}{ਉਹ ਸੰਖਿਆ ਜਿਸ ਨਾਲ ਕਿਸੇ ਆਕਾਰ ਨੂੰ ਵੱਡਾ ਕਰਨ ਵੇਲੇ ਹਰ ਲੰਬਾਈ ਨੂੰ ਗੁਣਾ ਕੀਤਾ ਜਾਂਦਾ ਹੈ।}
\dictentrytr{sector}{a slice of a circle between two radii}{ਸੈਕਟਰ}{ਦੋ ਅਰਧਵਿਆਸਾਂ ਦੇ ਵਿਚਕਾਰ ਚੱਕਰ ਦਾ ਟੁਕੜਾ}
\dictentrytr{segment}{the part of a circle cut off by a chord}{ਚੱਕਰ-ਖੰਡ}{ਚੱਕਰ ਦਾ ਉਹ ਹਿੱਸਾ ਜੋ ਇੱਕ ਜੀਵਾ ਦੁਆਰਾ ਕੱਟਿਆ ਗਿਆ ਹੋਵੇ}
\dictentrytr{semicircular}{shaped like half of a round object.}{ਅਰਧ-ਗੋਲਾਕਾਰ}{ਅੱਧੀ ਗੋਲ ਚੀਜ਼ ਵਾਂਗ ਆਕਾਰ ਵਾਲਾ।}
\dictentrytr{sequence}{a list of numbers that follows a rule}{ਲੜੀ}{ਸੰਖਿਆਵਾਂ ਦੀ ਸੂਚੀ ਜੋ ਇੱਕ ਨਿਯਮ ਦੀ ਪਾਲਣਾ ਕਰਦੀ ਹੈ}
\dictentrytr{set}{a collection of numbers or objects grouped together.}{ਸਮੂਹ}{ਸੰਖਿਆਵਾਂ ਜਾਂ ਵਸਤੂਆਂ ਦਾ ਸੰਗ੍ਰਹਿ ਜੋ ਇਕੱਠੇ ਕੀਤੀਆਂ ਗਈਆਂ ਹੋਣ।}
\dictentrytr{sign}{The plus or minus mark that tells whether a value is positive or negative.}{ਚਿੰਨ੍ਹ}{ਉਹ ਜੋੜ ਜਾਂ ਘਟਾਓ ਦਾ ਚਿੰਨ੍ਹ ਜੋ ਦੱਸਦਾ ਹੈ ਕਿ ਮੁੱਲ ਧਨਾਤਮਕ ਹੈ ਜਾਂ ਰਿਣਾਤਮਕ।}
\dictentrytr{similar}{the same shape but a different size}{ਸਮਰੂਪ}{ਇੱਕੋ ਸ਼ਕਲ ਪਰ ਵੱਖਰੇ ਆਕਾਰ ਵਾਲਾ}
\dictentrytr{simple interest}{interest calculated only on the original amount, so previous interest earns nothing.}{ਸਾਧਾਰਨ ਵਿਆਜ}{ਵਿਆਜ ਜੋ ਸਿਰਫ਼ ਅਸਲ ਰਾਸ਼ੀ ਉੱਤੇ ਲਗਾਇਆ ਜਾਂਦਾ ਹੈ, ਇਸ ਲਈ ਪਿਛਲਾ ਵਿਆਜ ਕੁਝ ਨਹੀਂ ਕਮਾਉਂਦਾ।}
\dictentrytr{simplest form}{a fraction written with the smallest possible numbers above and below the line.}{ਸਰਲਤਮ ਰੂਪ}{ਇੱਕ ਭਿੰਨ ਜੋ ਰੇਖਾ ਦੇ ਉੱਪਰ ਅਤੇ ਹੇਠਾਂ ਸਭ ਤੋਂ ਛੋਟੀਆਂ ਸੰਭਵ ਸੰਖਿਆਵਾਂ ਨਾਲ ਲਿਖੀ ਜਾਵੇ।}
\dictentrytr{simplify}{to write something in its smallest or plainest form without changing its value}{ਸਰਲ ਕਰਨਾ}{ਕਿਸੇ ਚੀਜ਼ ਨੂੰ ਉਸਦੇ ਸਭ ਤੋਂ ਛੋਟੇ ਜਾਂ ਸਾਦੇ ਰੂਪ ਵਿੱਚ ਲਿਖਣਾ, ਉਸਦਾ ਮੁੱਲ ਨਾ ਬਦਲਦੇ ਹੋਏ।}
\dictentrytr{simultaneous equations}{two or more equations with the same unknown values that are solved together.}{ਯੁਗਪਤ ਸਮੀਕਰਨ}{ਇੱਕੋ ਜਿਹੇ ਅਗਿਆਤ ਮੁੱਲਾਂ ਵਾਲੇ ਦੋ ਜਾਂ ਵੱਧ ਸਮੀਕਰਨ ਜੋ ਇਕੱਠੇ ਹੱਲ ਕੀਤੇ ਜਾਂਦੇ ਹਨ।}
\dictentrytr{sin}{the opposite side divided by the hypotenuse for an angle in a right triangle.}{ਸਾਈਨ}{ਸਮਕੋਣ ਤਿਕੋਣ ਵਿੱਚ ਕਿਸੇ ਕੋਣ ਲਈ ਸਾਹਮਣੇ ਵਾਲੀ ਭੁਜਾ ਨੂੰ ਕਰਨ ਨਾਲ ਭਾਗ ਕੀਤਾ ਜਾਂਦਾ ਹੈ।}
\dictentrytr{sine}{in a right-angled triangle, the side across from a chosen corner divided by the longest side.}{ਸਾਈਨ}{ਸਮਕੋਣ ਤਿਕੋਣ ਵਿੱਚ, ਕਿਸੇ ਚੁਣੇ ਕੋਣ ਦੇ ਸਾਹਮਣੇ ਵਾਲੀ ਭੁਜਾ ਨੂੰ ਸਭ ਤੋਂ ਲੰਬੀ ਭੁਜਾ ਨਾਲ ਭਾਗ ਕੀਤਾ ਜਾਂਦਾ ਹੈ।}
\dictentrytr{sine rule}{a formula connecting the sides of a triangle with the sin of the angle opposite each side.}{ਸਾਈਨ ਨਿਯਮ}{ਇੱਕ ਫਾਰਮੂਲਾ ਜੋ ਤਿਕੋਣ ਦੀਆਂ ਭੁਜਾਵਾਂ ਨੂੰ ਹਰ ਭੁਜਾ ਦੇ ਸਾਹਮਣੇ ਵਾਲੇ ਕੋਣ ਦੇ ਸਾਈਨ ਨਾਲ ਜੋੜਦਾ ਹੈ।}
\dictentrytr{solution}{a value or set of values that makes an equation true when used in place of the letter.}{ਹੱਲ}{ਉਹ ਮੁੱਲ ਜਾਂ ਮੁੱਲਾਂ ਦਾ ਸਮੂਹ ਜੋ ਅੱਖਰ ਦੀ ਥਾਂ ਵਰਤਣ ਨਾਲ ਸਮੀਕਰਨ ਨੂੰ ਸੱਚਾ ਬਣਾਉਂਦਾ ਹੈ।}
\dictentrytr{speed}{how far something travels in one unit of time}{ਗਤੀ}{ਕੋਈ ਚੀਜ਼ ਸਮੇਂ ਦੀ ਇੱਕ ਇਕਾਈ ਵਿੱਚ ਕਿੰਨੀ ਦੂਰੀ ਤੈਅ ਕਰਦੀ ਹੈ।}
\dictentrytr{square}{to work out a number times itself; also a flat shape with four equal sides and four equal corners}{ਵਰਗ}{ਕਿਸੇ ਸੰਖਿਆ ਨੂੰ ਆਪਣੇ ਨਾਲ ਗੁਣਾ ਕਰਨਾ; ਨਾਲ ਹੀ ਚਾਰ ਬਰਾਬਰ ਭੁਜਾਵਾਂ ਅਤੇ ਚਾਰ ਬਰਾਬਰ ਕੋਣਾਂ ਵਾਲਾ ਸਮਤਲ ਆਕਾਰ।}
\dictentrytr{square number}{a result of multiplying a whole number by itself, such as 1, 4, 9, 16.}{ਵਰਗ ਸੰਖਿਆ}{ਕਿਸੇ ਪੂਰਨ ਸੰਖਿਆ ਨੂੰ ਆਪਣੇ ਨਾਲ ਗੁਣਾ ਕਰਨ ਦਾ ਨਤੀਜਾ, ਜਿਵੇਂ 1, 4, 9, 16।}
\dictentrytr{substitute}{to put a number in place of a letter}{ਪ੍ਰਤਿਸਥਾਪਿਤ ਕਰਨਾ}{ਕਿਸੇ ਅੱਖਰ ਦੀ ਥਾਂ ਇੱਕ ਸੰਖਿਆ ਰੱਖਣੀ}
\dictentrytr{subtended}{formed when the two ends of a chord are joined to a point on the circumference.}{ਅੰਤਰਿਤ}{ਉਦੋਂ ਬਣਦਾ ਹੈ ਜਦੋਂ ਕਿਸੇ ਜੀਵਾ ਦੇ ਦੋਵੇਂ ਸਿਰੇ ਪਰਿਧੀ ਉੱਤੇ ਸਥਿਤ ਕਿਸੇ ਬਿੰਦੂ ਨਾਲ ਜੋੜੇ ਜਾਂਦੇ ਹਨ।}
\dictentrytr{subtraction}{the process of taking one amount away from another.}{ਘਟਾਓ}{ਇੱਕ ਮਾਤਰਾ ਵਿੱਚੋਂ ਦੂਜੀ ਮਾਤਰਾ ਨੂੰ ਘਟਾਉਣ ਦੀ ਪ੍ਰਕਿਰਿਆ।}
\dictentrytr{sum}{the answer when numbers are added together}{ਜੋੜਫਲ}{ਜਦੋਂ ਸੰਖਿਆਵਾਂ ਨੂੰ ਜੋੜਿਆ ਜਾਂਦਾ ਹੈ ਤਾਂ ਮਿਲਣ ਵਾਲਾ ਉੱਤਰ।}
\dictentrytr{surd}{a root that cannot be written exactly as a fraction, so it is left in root form}{ਕਰਣੀ}{ਇੱਕ ਅਜਿਹਾ ਮੂਲ ਜੋ ਭਿੰਨ ਦੇ ਰੂਪ ਵਿੱਚ ਠੀਕ-ਠੀਕ ਨਹੀਂ ਲਿਖਿਆ ਜਾ ਸਕਦਾ, ਇਸ ਲਈ ਉਹ ਮੂਲ ਰੂਪ ਵਿੱਚ ਰਹਿੰਦਾ ਹੈ।}
\dictentrytr{symmetric}{balanced so that one side is a mirror image of the other.}{ਸਮਮਿਤ}{ਇਸ ਤਰ੍ਹਾਂ ਸੰਤੁਲਿਤ ਕਿ ਇੱਕ ਪਾਸਾ ਦੂਜੇ ਪਾਸੇ ਦਾ ਦਰਪਣ ਪ੍ਰਤੀਬਿੰਬ ਹੋਵੇ।}
\dictentrytr{tan}{the opposite side divided by the adjacent side for an angle in a right triangle.}{ਟੈਨ}{ਸਮਕੋਣ ਤਿਕੋਣ ਵਿੱਚ ਕਿਸੇ ਕੋਣ ਲਈ ਸਾਹਮਣੇ ਵਾਲੀ ਭੁਜਾ ਨੂੰ ਲਾਗਲੀ ਭੁਜਾ ਨਾਲ ਭਾਗ ਕੀਤਾ ਜਾਂਦਾ ਹੈ।}
\dictentrytr{tangent}{a line that touches a curve at exactly one point}{ਸਪਰਸ਼ ਰੇਖਾ}{ਇੱਕ ਰੇਖਾ ਜੋ ਕਿਸੇ ਵਕਰ ਨੂੰ ਬਿਲਕੁਲ ਇੱਕ ਬਿੰਦੂ ਉੱਤੇ ਛੋਹਂਦੀ ਹੈ}
\dictentrytr{ten-thousandths}{the value of one part when a whole is divided into ten thousand equal parts.}{ਦਸ-ਹਜ਼ਾਰਵੇਂ}{ਇੱਕ ਪੂਰੇ ਨੂੰ ਦਸ ਹਜ਼ਾਰ ਬਰਾਬਰ ਭਾਗਾਂ ਵਿੱਚ ਵੰਡਣ ਨਾਲ ਇੱਕ ਭਾਗ ਦਾ ਮੁੱਲ।}
\dictentrytr{tenths}{the value of one part when a whole is divided into ten equal parts.}{ਦਸਵੇਂ}{ਇੱਕ ਪੂਰੇ ਨੂੰ ਦਸ ਬਰਾਬਰ ਭਾਗਾਂ ਵਿੱਚ ਵੰਡਣ ਨਾਲ ਇੱਕ ਭਾਗ ਦਾ ਮੁੱਲ।}
\dictentrytr{term}{one part of an expression or one number in a sequence}{ਪਦ}{ਕਿਸੇ ਵਿਅੰਜਕ ਦਾ ਇੱਕ ਹਿੱਸਾ ਜਾਂ ਲੜੀ ਵਿੱਚ ਇੱਕ ਸੰਖਿਆ}
\dictentrytr{terminate}{for a decimal to end after a fixed number of digits.}{ਸਮਾਪਤ ਹੋਣਾ}{ਦਸ਼ਮਲਵ ਦਾ ਨਿਸ਼ਚਿਤ ਅੰਕਾਂ ਤੋਂ ਬਾਅਦ ਖਤਮ ਹੋਣਾ।}
\dictentrytr{theorem}{a mathematical statement that can be proved to be true.}{ਪ੍ਰਮੇਯ}{ਇੱਕ ਗਣਿਤ ਦਾ ਕਥਨ ਜੋ ਸੱਚ ਸਾਬਤ ਕੀਤਾ ਜਾ ਸਕੇ।}
\dictentrytr{thousandths}{the value of one part when a whole is divided into one thousand equal parts.}{ਹਜ਼ਾਰਵੇਂ}{ਇੱਕ ਪੂਰੇ ਨੂੰ ਇੱਕ ਹਜ਼ਾਰ ਬਰਾਬਰ ਭਾਗਾਂ ਵਿੱਚ ਵੰਡਣ ਨਾਲ ਇੱਕ ਭਾਗ ਦਾ ਮੁੱਲ।}
\dictentrytr{times table}{a list of numbers made by repeatedly adding the same number, such as 3, 6, 9, 12.}{ਪਹਾੜਾ}{ਇੱਕੋ ਸੰਖਿਆ ਨੂੰ ਵਾਰ-ਵਾਰ ਜੋੜ ਕੇ ਬਣਾਈ ਗਈ ਸੰਖਿਆਵਾਂ ਦੀ ਸੂਚੀ, ਜਿਵੇਂ 3, 6, 9, 12।}
\dictentrytr{tip-to-tail}{A method for adding vectors by placing the point of one at the start of the next.}{ਸਿਰੇ-ਤੋਂ-ਪੂਛ}{ਵੈਕਟਰਾਂ ਨੂੰ ਜੋੜਨ ਦਾ ਤਰੀਕਾ ਜਿਸ ਵਿੱਚ ਇੱਕ ਦੇ ਸਿਰੇ ਨੂੰ ਅਗਲੇ ਦੇ ਸ਼ੁਰੂ ਉੱਤੇ ਰੱਖਿਆ ਜਾਂਦਾ ਹੈ।}
\dictentrytr{total}{the whole amount found by adding all the parts together.}{ਕੁੱਲ}{ਸਾਰੇ ਭਾਗਾਂ ਨੂੰ ਜੋੜ ਕੇ ਮਿਲਣ ਵਾਲੀ ਪੂਰੀ ਮਾਤਰਾ।}
\dictentrytr{transformation}{a change to a shape's position, size or orientation}{ਰੂਪਾਂਤਰਨ}{ਕਿਸੇ ਆਕਾਰ ਦੀ ਸਥਿਤੀ, ਆਕਾਰ ਜਾਂ ਦਿਸ਼ਾ ਵਿੱਚ ਤਬਦੀਲੀ}
\dictentrytr{translation}{moving every point by the same vector, without turning or resizing it}{ਸਥਾਨਾਂਤਰਨ}{ਹਰ ਬਿੰਦੂ ਨੂੰ ਇੱਕੋ ਵੈਕਟਰ ਨਾਲ ਹਿਲਾਉਣਾ, ਬਿਨਾਂ ਘੁਮਾਏ ਜਾਂ ਆਕਾਰ ਬਦਲੇ।}
\dictentrytr{triangle}{a flat shape with three straight sides}{ਤਿਕੋਣ}{ਤਿੰਨ ਸਿੱਧੀਆਂ ਭੁਜਾਵਾਂ ਵਾਲਾ ਸਮਤਲ ਆਕਾਰ।}
\dictentrytr{triangular number}{a number that can be shown as dots arranged in a triangle}{ਤਿਕੋਣੀ ਸੰਖਿਆ}{ਉਹ ਸੰਖਿਆ ਜੋ ਬਿੰਦੂਆਂ ਨੂੰ ਤਿਕੋਣ ਵਿੱਚ ਲਗਾ ਕੇ ਦਿਖਾਈ ਜਾ ਸਕੇ।}
\dictentrytr{trigonometry}{the mathematics of relationships between angles and side lengths in triangles.}{ਤ੍ਰਿਕੋਣਮਿਤੀ}{ਤਿਕੋਣਾਂ ਵਿੱਚ ਕੋਣਾਂ ਅਤੇ ਭੁਜਾਵਾਂ ਦੀਆਂ ਲੰਬਾਈਆਂ ਵਿਚਲੇ ਸਬੰਧਾਂ ਦਾ ਗਣਿਤ।}
\dictentrytr{uniform}{the same at all points; not changing.}{ਇੱਕਸਮਾਨ}{ਹਰ ਬਿੰਦੂ ਉੱਤੇ ਇੱਕੋ ਜਿਹਾ; ਨਾ ਬਦਲਣ ਵਾਲਾ।}
\dictentrytr{unique}{being the only one; here, no other value is tied with the mode.}{ਇਕਲੌਤਾ}{ਸਿਰਫ਼ ਇੱਕ ਹੋਣਾ; ਇੱਥੇ, ਕੋਈ ਹੋਰ ਮੁੱਲ ਬਹੁਲਕ ਨਾਲ ਬਰਾਬਰ ਨਹੀਂ ਹੈ।}
\dictentrytr{unit fraction}{a fraction whose top number is one.}{ਇਕਾਈ ਭਿੰਨ}{ਇੱਕ ਭਿੰਨ ਜਿਸਦੀ ਉੱਪਰਲੀ ਸੰਖਿਆ ਇੱਕ ਹੁੰਦੀ ਹੈ।}
\dictentrytr{unit price}{the cost for one single gram, litre, item, or other measure}{ਇਕਾਈ ਮੁੱਲ}{ਇੱਕ ਗ੍ਰਾਮ, ਲਿਟਰ, ਵਸਤੂ ਜਾਂ ਕਿਸੇ ਹੋਰ ਮਾਪ ਦੀ ਕੀਮਤ।}
\dictentrytr{unitary}{a method that first finds the value of one part or one item}{ਇਕਾਈ ਵਿਧੀ}{ਇੱਕ ਵਿਧੀ ਜੋ ਪਹਿਲਾਂ ਇੱਕ ਹਿੱਸੇ ਜਾਂ ਇੱਕ ਵਸਤੂ ਦਾ ਮੁੱਲ ਕੱਢਦੀ ਹੈ।}
\dictentrytr{upper bound}{a number that is greater than or equal to every possible value in a set.}{ਉੱਪਰਲੀ ਸੀਮਾ}{ਉਹ ਸੰਖਿਆ ਜੋ ਕਿਸੇ ਸਮੂਹ ਵਿੱਚ ਹਰ ਸੰਭਵ ਮੁੱਲ ਤੋਂ ਵੱਡੀ ਜਾਂ ਬਰਾਬਰ ਹੋਵੇ।}
\dictentrytr{value}{the number that a letter or symbol stands for}{ਮੁੱਲ}{ਉਹ ਸੰਖਿਆ ਜੋ ਕਿਸੇ ਅੱਖਰ ਜਾਂ ਚਿੰਨ੍ਹ ਦੁਆਰਾ ਦਰਸਾਈ ਜਾਂਦੀ ਹੈ।}
\dictentrytr{variable}{a letter standing for a quantity that can change}{ਚਲ}{ਇੱਕ ਅੱਖਰ ਜੋ ਕਿਸੇ ਅਜਿਹੀ ਰਾਸ਼ੀ ਦੀ ਨੁਮਾਇੰਦਗੀ ਕਰਦਾ ਹੈ ਜੋ ਬਦਲ ਸਕਦੀ ਹੈ}
\dictentrytr{vector}{a quantity with both size and direction}{ਸਦਿਸ਼}{ਇੱਕ ਰਾਸ਼ੀ ਜਿਸ ਵਿੱਚ ਆਕਾਰ ਅਤੇ ਦਿਸ਼ਾ ਦੋਵੇਂ ਹੁੰਦੇ ਹਨ}
\dictentrytr{vertex}{a corner where two or more edges meet}{ਸਿਖਰ}{ਇੱਕ ਕੋਨਾ ਜਿੱਥੇ ਦੋ ਜਾਂ ਵੱਧ ਕਿਨਾਰੇ ਮਿਲਦੇ ਹਨ}
\dictentrytr{vertical}{going straight up and down}{ਲੰਬਕਾਰੀ}{ਸਿੱਧਾ ਉੱਪਰ ਅਤੇ ਹੇਠਾਂ ਜਾਂਦਾ ਹੋਇਆ।}
\dictentrytr{vertically opposite}{describing two angles, the same size, formed across from each other when two lines cross}{ਸਿਖਰ-ਸਨਮੁਖ}{ਦੋ ਕੋਣਾਂ ਦਾ ਵਰਣਨ ਕਰਦਾ ਹੈ ਜੋ ਇੱਕੋ ਮਾਪ ਦੇ ਹੁੰਦੇ ਹਨ ਅਤੇ ਦੋ ਰੇਖਾਵਾਂ ਦੇ ਇੱਕ-ਦੂਜੇ ਨੂੰ ਕੱਟਣ ਨਾਲ ਇੱਕ-ਦੂਜੇ ਦੇ ਸਾਹਮਣੇ ਬਣਦੇ ਹਨ।}
\dictentrytr{volume}{the amount of space inside a solid shape}{ਆਇਤਨ}{ਕਿਸੇ ਠੋਸ ਆਕਾਰ ਦੇ ਅੰਦਰ ਥਾਂ ਦੀ ਮਾਤਰਾ}
\dictentrytr{whole number}{a number that is not negative and has no part after a decimal point, such as 0, 1, 2 or 3.}{ਪੂਰਨ ਸੰਖਿਆ}{ਇੱਕ ਸੰਖਿਆ ਜੋ ਰਿਣਾਤਮਕ ਨਹੀਂ ਹੁੰਦੀ ਅਤੇ ਦਸ਼ਮਲਵ ਬਿੰਦੂ ਤੋਂ ਬਾਅਦ ਉਸਦਾ ਕੋਈ ਹਿੱਸਾ ਨਹੀਂ ਹੁੰਦਾ, ਜਿਵੇਂ 0, 1, 2 ਜਾਂ 3।}
\dictentrytr{whole numbers}{the numbers zero, one, two, three, and so on, with no fraction or decimal part.}{ਪੂਰਨ ਸੰਖਿਆਵਾਂ}{ਸਿਫ਼ਰ, ਇੱਕ, ਦੋ, ਤਿੰਨ, ਆਦਿ ਸੰਖਿਆਵਾਂ, ਜਿਨ੍ਹਾਂ ਵਿੱਚ ਕੋਈ ਭਿੰਨ ਜਾਂ ਦਸ਼ਮਲਵ ਭਾਗ ਨਹੀਂ ਹੁੰਦਾ।}
\dictentrytr{yield}{The money earned on an amount lent or invested, usually shown as a percentage.}{ਪ੍ਰਤਿਫਲ}{ਉਧਾਰ ਦਿੱਤੀ ਜਾਂ ਨਿਵੇਸ਼ ਕੀਤੀ ਰਕਮ 'ਤੇ ਕਮਾਇਆ ਗਿਆ ਪੈਸਾ, ਜੋ ਆਮ ਤੌਰ 'ਤੇ ਪ੍ਰਤੀਸ਼ਤ ਵਜੋਂ ਦਰਸਾਇਆ ਜਾਂਦਾ ਹੈ।}
\dictentrytr{yield curve}{a line showing how yields change for different times to maturity.}{ਪ੍ਰਤਿਫਲ ਵਕਰ}{ਇੱਕ ਰੇਖਾ ਜੋ ਦਿਖਾਉਂਦੀ ਹੈ ਕਿ ਪਰਿਪੱਕਤਾ ਦੇ ਵੱਖ-ਵੱਖ ਸਮਿਆਂ ਲਈ ਪ੍ਰਤਿਫਲ ਕਿਵੇਂ ਬਦਲਦਾ ਹੈ।}

\end{document}
